\subsection*{Textbook Formal Inventory}
This subsection records the exact formal material in \file{HAETAE_Keccak1600.ec}.  It is generated from the proof-surface EasyCrypt file, so the line numbers and statements are meant to be read next to the checked source.

\noindent\textbf{Chapter role.} Keccak-f1600 state and bit-level round equations.

\noindent\textbf{Inventory size.} 2 context declarations, 37 definitions/game objects, 50 lemmas or relational theorems, and 0 explicit Hoare/Phoare judgments were found.

\subsubsection*{Theory Context, Imports, and Quantified Modules}
\paragraph{Context 1 (line 1)}
\noindent\textbf{EasyCrypt name.}
\begin{lstlisting}[language=EasyCrypt,basicstyle=\ttfamily\scriptsize]
imported theories
\end{lstlisting}
\noindent\textbf{Checked EasyCrypt statement.}
\begin{lstlisting}[language=EasyCrypt,basicstyle=\ttfamily\scriptsize]
require import AllCore IntDiv List.
\end{lstlisting}
\noindent\textbf{Formal mathematical reading.}
Mathematical context. This line imports or specializes earlier theories, making their carrier sets, functions, modules, and theorems available to the current file.

\noindent\textbf{Role in the proof.}
Use in this chapter. It fixes the proof context, imported mathematical language, or quantified module parameters for the following statements.

\paragraph{Context 2 (line 3)}
\noindent\textbf{EasyCrypt name.}
\begin{lstlisting}[language=EasyCrypt,basicstyle=\ttfamily\scriptsize]
HAETAE_Keccak1600
\end{lstlisting}
\noindent\textbf{Checked EasyCrypt statement.}
\begin{lstlisting}[language=EasyCrypt,basicstyle=\ttfamily\scriptsize]
theory HAETAE_Keccak1600.
\end{lstlisting}
\noindent\textbf{Formal mathematical reading.}
Mathematical context. This begins a namespace for the formal development in this file.

\noindent\textbf{Role in the proof.}
Use in this chapter. It fixes the proof context, imported mathematical language, or quantified module parameters for the following statements.

\subsubsection*{Definitions, Carrier Sets, Operations, Games, and Procedures}
\begin{formaldefinition}[EasyCrypt line 5]
\noindent\textbf{EasyCrypt name.}
\begin{lstlisting}[language=EasyCrypt,basicstyle=\ttfamily\scriptsize]
keccak_byte
\end{lstlisting}
\noindent\textbf{Checked EasyCrypt statement.}
\begin{lstlisting}[language=EasyCrypt,basicstyle=\ttfamily\scriptsize]
type keccak_byte = int.
\end{lstlisting}
\noindent\textbf{Formal mathematical reading.}
Mathematical definition. This is a definitional type abbreviation.  The exact displayed EasyCrypt statement gives the right-hand side; later proof steps may unfold it without adding an assumption.

\noindent\textbf{Role in the proof.}
Use in this chapter. It is part of the exact formal vocabulary used by subsequent lemmas and games.

\end{formaldefinition}

\begin{formaldefinition}[EasyCrypt line 6]
\noindent\textbf{EasyCrypt name.}
\begin{lstlisting}[language=EasyCrypt,basicstyle=\ttfamily\scriptsize]
keccak_lane
\end{lstlisting}
\noindent\textbf{Checked EasyCrypt statement.}
\begin{lstlisting}[language=EasyCrypt,basicstyle=\ttfamily\scriptsize]
type keccak_lane = bool list.
\end{lstlisting}
\noindent\textbf{Formal mathematical reading.}
Mathematical definition. This is a definitional type abbreviation.  The exact displayed EasyCrypt statement gives the right-hand side; later proof steps may unfold it without adding an assumption.

\noindent\textbf{Role in the proof.}
Use in this chapter. It is part of the exact formal vocabulary used by subsequent lemmas and games.

\end{formaldefinition}

\begin{formaldefinition}[EasyCrypt line 7]
\noindent\textbf{EasyCrypt name.}
\begin{lstlisting}[language=EasyCrypt,basicstyle=\ttfamily\scriptsize]
keccak_state
\end{lstlisting}
\noindent\textbf{Checked EasyCrypt statement.}
\begin{lstlisting}[language=EasyCrypt,basicstyle=\ttfamily\scriptsize]
type keccak_state = keccak_lane list.
\end{lstlisting}
\noindent\textbf{Formal mathematical reading.}
Mathematical definition. This is a definitional type abbreviation.  The exact displayed EasyCrypt statement gives the right-hand side; later proof steps may unfold it without adding an assumption.

\noindent\textbf{Role in the proof.}
Use in this chapter. It is part of the exact formal vocabulary used by subsequent lemmas and games.

\end{formaldefinition}

\begin{formaldefinition}[EasyCrypt line 9]
\noindent\textbf{EasyCrypt name.}
\begin{lstlisting}[language=EasyCrypt,basicstyle=\ttfamily\scriptsize]
keccak_lane_bits
\end{lstlisting}
\noindent\textbf{Checked EasyCrypt statement.}
\begin{lstlisting}[language=EasyCrypt,basicstyle=\ttfamily\scriptsize]
op keccak_lane_bits : int = 64.
\end{lstlisting}
\noindent\textbf{Formal mathematical reading.}
Mathematical definition. This is a total EasyCrypt operation.  The exact displayed statement gives the domain, codomain, and defining equation when present.

\noindent\textbf{Role in the proof.}
Use in this chapter. It is part of the exact formal vocabulary used by subsequent lemmas and games.

\end{formaldefinition}

\begin{formaldefinition}[EasyCrypt line 10]
\noindent\textbf{EasyCrypt name.}
\begin{lstlisting}[language=EasyCrypt,basicstyle=\ttfamily\scriptsize]
keccak_state_lanes
\end{lstlisting}
\noindent\textbf{Checked EasyCrypt statement.}
\begin{lstlisting}[language=EasyCrypt,basicstyle=\ttfamily\scriptsize]
op keccak_state_lanes : int = 25.
\end{lstlisting}
\noindent\textbf{Formal mathematical reading.}
Mathematical definition. This is a total EasyCrypt operation.  The exact displayed statement gives the domain, codomain, and defining equation when present.

\noindent\textbf{Role in the proof.}
Use in this chapter. It is part of the exact formal vocabulary used by subsequent lemmas and games.

\end{formaldefinition}

\begin{formaldefinition}[EasyCrypt line 11]
\noindent\textbf{EasyCrypt name.}
\begin{lstlisting}[language=EasyCrypt,basicstyle=\ttfamily\scriptsize]
keccak_state_bytes
\end{lstlisting}
\noindent\textbf{Checked EasyCrypt statement.}
\begin{lstlisting}[language=EasyCrypt,basicstyle=\ttfamily\scriptsize]
op keccak_state_bytes : int = 200.
\end{lstlisting}
\noindent\textbf{Formal mathematical reading.}
Mathematical definition. This is a total EasyCrypt operation.  The exact displayed statement gives the domain, codomain, and defining equation when present.

\noindent\textbf{Role in the proof.}
Use in this chapter. It is part of the exact formal vocabulary used by subsequent lemmas and games.

\end{formaldefinition}

\begin{formaldefinition}[EasyCrypt line 13]
\noindent\textbf{EasyCrypt name.}
\begin{lstlisting}[language=EasyCrypt,basicstyle=\ttfamily\scriptsize]
keccak_byte_norm
\end{lstlisting}
\noindent\textbf{Checked EasyCrypt statement.}
\begin{lstlisting}[language=EasyCrypt,basicstyle=\ttfamily\scriptsize]
op keccak_byte_norm (b : keccak_byte) : keccak_byte = b %% 256.
\end{lstlisting}
\noindent\textbf{Formal mathematical reading.}
Mathematical definition. This is a total EasyCrypt operation.  The exact displayed statement gives the domain, codomain, and defining equation when present.

\noindent\textbf{Role in the proof.}
Use in this chapter. It is part of the exact formal vocabulary used by subsequent lemmas and games.

\end{formaldefinition}

\begin{formaldefinition}[EasyCrypt line 14]
\noindent\textbf{EasyCrypt name.}
\begin{lstlisting}[language=EasyCrypt,basicstyle=\ttfamily\scriptsize]
keccak_byte_bit
\end{lstlisting}
\noindent\textbf{Checked EasyCrypt statement.}
\begin{lstlisting}[language=EasyCrypt,basicstyle=\ttfamily\scriptsize]
op keccak_byte_bit (b : keccak_byte) (i : int) : bool =
  ((keccak_byte_norm b) %/ (2 ^ i)) %% 2 <> 0.
\end{lstlisting}
\noindent\textbf{Formal mathematical reading.}
Mathematical definition. This is a Boolean predicate.  The exact displayed EasyCrypt statement gives its arguments; mathematically it is an event, invariant, support condition, or well-formedness predicate over those arguments.

\noindent\textbf{Role in the proof.}
Use in this chapter. It is part of the exact formal vocabulary used by subsequent lemmas and games.

\end{formaldefinition}

\begin{formaldefinition}[EasyCrypt line 16]
\noindent\textbf{EasyCrypt name.}
\begin{lstlisting}[language=EasyCrypt,basicstyle=\ttfamily\scriptsize]
keccak_word_norm
\end{lstlisting}
\noindent\textbf{Checked EasyCrypt statement.}
\begin{lstlisting}[language=EasyCrypt,basicstyle=\ttfamily\scriptsize]
op keccak_word_norm (w : int) : int = w %% (2 ^ keccak_lane_bits).
\end{lstlisting}
\noindent\textbf{Formal mathematical reading.}
Mathematical definition. This is a total EasyCrypt operation.  The exact displayed statement gives the domain, codomain, and defining equation when present.

\noindent\textbf{Role in the proof.}
Use in this chapter. It is part of the exact formal vocabulary used by subsequent lemmas and games.

\end{formaldefinition}

\begin{formaldefinition}[EasyCrypt line 17]
\noindent\textbf{EasyCrypt name.}
\begin{lstlisting}[language=EasyCrypt,basicstyle=\ttfamily\scriptsize]
keccak_word_bit
\end{lstlisting}
\noindent\textbf{Checked EasyCrypt statement.}
\begin{lstlisting}[language=EasyCrypt,basicstyle=\ttfamily\scriptsize]
op keccak_word_bit (w : int) (i : int) : bool =
  ((keccak_word_norm w) %/ (2 ^ i)) %% 2 <> 0.
\end{lstlisting}
\noindent\textbf{Formal mathematical reading.}
Mathematical definition. This is a Boolean predicate.  The exact displayed EasyCrypt statement gives its arguments; mathematically it is an event, invariant, support condition, or well-formedness predicate over those arguments.

\noindent\textbf{Role in the proof.}
Use in this chapter. It is part of the exact formal vocabulary used by subsequent lemmas and games.

\end{formaldefinition}

\begin{formaldefinition}[EasyCrypt line 20]
\noindent\textbf{EasyCrypt name.}
\begin{lstlisting}[language=EasyCrypt,basicstyle=\ttfamily\scriptsize]
keccak_lane_zero
\end{lstlisting}
\noindent\textbf{Checked EasyCrypt statement.}
\begin{lstlisting}[language=EasyCrypt,basicstyle=\ttfamily\scriptsize]
op keccak_lane_zero : keccak_lane = nseq keccak_lane_bits false.
\end{lstlisting}
\noindent\textbf{Formal mathematical reading.}
Mathematical definition. This is a total EasyCrypt operation.  The exact displayed statement gives the domain, codomain, and defining equation when present.

\noindent\textbf{Role in the proof.}
Use in this chapter. It is part of the exact formal vocabulary used by subsequent lemmas and games.

\end{formaldefinition}

\begin{formaldefinition}[EasyCrypt line 21]
\noindent\textbf{EasyCrypt name.}
\begin{lstlisting}[language=EasyCrypt,basicstyle=\ttfamily\scriptsize]
keccak_lane_bit
\end{lstlisting}
\noindent\textbf{Checked EasyCrypt statement.}
\begin{lstlisting}[language=EasyCrypt,basicstyle=\ttfamily\scriptsize]
op keccak_lane_bit (lane : keccak_lane) (i : int) : bool =
  nth false lane i.
\end{lstlisting}
\noindent\textbf{Formal mathematical reading.}
Mathematical definition. This is a Boolean predicate.  The exact displayed EasyCrypt statement gives its arguments; mathematically it is an event, invariant, support condition, or well-formedness predicate over those arguments.

\noindent\textbf{Role in the proof.}
Use in this chapter. It is part of the exact formal vocabulary used by subsequent lemmas and games.

\end{formaldefinition}

\begin{formaldefinition}[EasyCrypt line 23]
\noindent\textbf{EasyCrypt name.}
\begin{lstlisting}[language=EasyCrypt,basicstyle=\ttfamily\scriptsize]
keccak_lane_wf
\end{lstlisting}
\noindent\textbf{Checked EasyCrypt statement.}
\begin{lstlisting}[language=EasyCrypt,basicstyle=\ttfamily\scriptsize]
op keccak_lane_wf (lane : keccak_lane) : bool =
  size lane = keccak_lane_bits.
\end{lstlisting}
\noindent\textbf{Formal mathematical reading.}
Mathematical definition. This is a Boolean predicate.  The exact displayed EasyCrypt statement gives its arguments; mathematically it is an event, invariant, support condition, or well-formedness predicate over those arguments.

\noindent\textbf{Role in the proof.}
Use in this chapter. It names a well-formedness, support, or validity predicate needed by later preservation lemmas.

\end{formaldefinition}

\begin{formaldefinition}[EasyCrypt line 26]
\noindent\textbf{EasyCrypt name.}
\begin{lstlisting}[language=EasyCrypt,basicstyle=\ttfamily\scriptsize]
keccak_lane_xor
\end{lstlisting}
\noindent\textbf{Checked EasyCrypt statement.}
\begin{lstlisting}[language=EasyCrypt,basicstyle=\ttfamily\scriptsize]
op keccak_lane_xor (a b : keccak_lane) : keccak_lane =
  mkseq
    (fun i => keccak_lane_bit a i <> keccak_lane_bit b i)
    keccak_lane_bits.
\end{lstlisting}
\noindent\textbf{Formal mathematical reading.}
Mathematical definition. This is a total EasyCrypt operation.  The exact displayed statement gives the domain, codomain, and defining equation when present.

\noindent\textbf{Role in the proof.}
Use in this chapter. It is part of the exact formal vocabulary used by subsequent lemmas and games.

\end{formaldefinition}

\begin{formaldefinition}[EasyCrypt line 30]
\noindent\textbf{EasyCrypt name.}
\begin{lstlisting}[language=EasyCrypt,basicstyle=\ttfamily\scriptsize]
keccak_lane_and
\end{lstlisting}
\noindent\textbf{Checked EasyCrypt statement.}
\begin{lstlisting}[language=EasyCrypt,basicstyle=\ttfamily\scriptsize]
op keccak_lane_and (a b : keccak_lane) : keccak_lane =
  mkseq
    (fun i => keccak_lane_bit a i /\ keccak_lane_bit b i)
    keccak_lane_bits.
\end{lstlisting}
\noindent\textbf{Formal mathematical reading.}
Mathematical definition. This is a total EasyCrypt operation.  The exact displayed statement gives the domain, codomain, and defining equation when present.

\noindent\textbf{Role in the proof.}
Use in this chapter. It is part of the exact formal vocabulary used by subsequent lemmas and games.

\end{formaldefinition}

\begin{formaldefinition}[EasyCrypt line 34]
\noindent\textbf{EasyCrypt name.}
\begin{lstlisting}[language=EasyCrypt,basicstyle=\ttfamily\scriptsize]
keccak_lane_not
\end{lstlisting}
\noindent\textbf{Checked EasyCrypt statement.}
\begin{lstlisting}[language=EasyCrypt,basicstyle=\ttfamily\scriptsize]
op keccak_lane_not (a : keccak_lane) : keccak_lane =
  mkseq
    (fun i => if keccak_lane_bit a i then false else true)
    keccak_lane_bits.
\end{lstlisting}
\noindent\textbf{Formal mathematical reading.}
Mathematical definition. This is a total EasyCrypt operation.  The exact displayed statement gives the domain, codomain, and defining equation when present.

\noindent\textbf{Role in the proof.}
Use in this chapter. It is part of the exact formal vocabulary used by subsequent lemmas and games.

\end{formaldefinition}

\begin{formaldefinition}[EasyCrypt line 38]
\noindent\textbf{EasyCrypt name.}
\begin{lstlisting}[language=EasyCrypt,basicstyle=\ttfamily\scriptsize]
keccak_lane_rotl
\end{lstlisting}
\noindent\textbf{Checked EasyCrypt statement.}
\begin{lstlisting}[language=EasyCrypt,basicstyle=\ttfamily\scriptsize]
op keccak_lane_rotl (a : keccak_lane) (offset : int) : keccak_lane =
  mkseq
    (fun i => keccak_lane_bit a ((i - offset) %% keccak_lane_bits))
    keccak_lane_bits.
\end{lstlisting}
\noindent\textbf{Formal mathematical reading.}
Mathematical definition. This is a total EasyCrypt operation.  The exact displayed statement gives the domain, codomain, and defining equation when present.

\noindent\textbf{Role in the proof.}
Use in this chapter. It is part of the exact formal vocabulary used by subsequent lemmas and games.

\end{formaldefinition}

\begin{formaldefinition}[EasyCrypt line 43]
\noindent\textbf{EasyCrypt name.}
\begin{lstlisting}[language=EasyCrypt,basicstyle=\ttfamily\scriptsize]
keccak_lane_of_int
\end{lstlisting}
\noindent\textbf{Checked EasyCrypt statement.}
\begin{lstlisting}[language=EasyCrypt,basicstyle=\ttfamily\scriptsize]
op keccak_lane_of_int (x : int) : keccak_lane =
  mkseq (fun i => keccak_word_bit x i) keccak_lane_bits.
\end{lstlisting}
\noindent\textbf{Formal mathematical reading.}
Mathematical definition. This is a total EasyCrypt operation.  The exact displayed statement gives the domain, codomain, and defining equation when present.

\noindent\textbf{Role in the proof.}
Use in this chapter. It is part of the exact formal vocabulary used by subsequent lemmas and games.

\end{formaldefinition}

\begin{formaldefinition}[EasyCrypt line 46]
\noindent\textbf{EasyCrypt name.}
\begin{lstlisting}[language=EasyCrypt,basicstyle=\ttfamily\scriptsize]
keccak_lane_to_byte
\end{lstlisting}
\noindent\textbf{Checked EasyCrypt statement.}
\begin{lstlisting}[language=EasyCrypt,basicstyle=\ttfamily\scriptsize]
op keccak_lane_to_byte (lane : keccak_lane) (byte_index : int) : keccak_byte =
  sumz
    (map
      (fun j =>
         if keccak_lane_bit lane (8 * byte_index + j)
         then 2 ^ j
         else 0)
      (range 0 8)).
\end{lstlisting}
\noindent\textbf{Formal mathematical reading.}
Mathematical definition. This is a total EasyCrypt operation.  The exact displayed statement gives the domain, codomain, and defining equation when present.

\noindent\textbf{Role in the proof.}
Use in this chapter. It is part of the exact formal vocabulary used by subsequent lemmas and games.

\end{formaldefinition}

\begin{formaldefinition}[EasyCrypt line 55]
\noindent\textbf{EasyCrypt name.}
\begin{lstlisting}[language=EasyCrypt,basicstyle=\ttfamily\scriptsize]
keccak_bytes_to_lane
\end{lstlisting}
\noindent\textbf{Checked EasyCrypt statement.}
\begin{lstlisting}[language=EasyCrypt,basicstyle=\ttfamily\scriptsize]
op keccak_bytes_to_lane (bs : keccak_byte list) (lane_index : int) :
  keccak_lane =
  mkseq
    (fun bit =>
       keccak_byte_bit
         (nth 0 bs (8 * lane_index + bit %/ 8))
         (bit %% 8))
    keccak_lane_bits.
\end{lstlisting}
\noindent\textbf{Formal mathematical reading.}
Mathematical definition. This is a total EasyCrypt operation.  The exact displayed statement gives the domain, codomain, and defining equation when present.

\noindent\textbf{Role in the proof.}
Use in this chapter. It is part of the exact formal vocabulary used by subsequent lemmas and games.

\end{formaldefinition}

\begin{formaldefinition}[EasyCrypt line 64]
\noindent\textbf{EasyCrypt name.}
\begin{lstlisting}[language=EasyCrypt,basicstyle=\ttfamily\scriptsize]
keccak_lanes_of_bytes
\end{lstlisting}
\noindent\textbf{Checked EasyCrypt statement.}
\begin{lstlisting}[language=EasyCrypt,basicstyle=\ttfamily\scriptsize]
op keccak_lanes_of_bytes (bs : keccak_byte list) : keccak_state =
  mkseq (fun lane => keccak_bytes_to_lane bs lane) keccak_state_lanes.
\end{lstlisting}
\noindent\textbf{Formal mathematical reading.}
Mathematical definition. This is a total EasyCrypt operation.  The exact displayed statement gives the domain, codomain, and defining equation when present.

\noindent\textbf{Role in the proof.}
Use in this chapter. It is part of the exact formal vocabulary used by subsequent lemmas and games.

\end{formaldefinition}

\begin{formaldefinition}[EasyCrypt line 67]
\noindent\textbf{EasyCrypt name.}
\begin{lstlisting}[language=EasyCrypt,basicstyle=\ttfamily\scriptsize]
keccak_bytes_of_lanes
\end{lstlisting}
\noindent\textbf{Checked EasyCrypt statement.}
\begin{lstlisting}[language=EasyCrypt,basicstyle=\ttfamily\scriptsize]
op keccak_bytes_of_lanes (st : keccak_state) : keccak_byte list =
  mkseq
    (fun i =>
       keccak_lane_to_byte
         (nth keccak_lane_zero st (i %/ 8))
         (i %% 8))
    keccak_state_bytes.
\end{lstlisting}
\noindent\textbf{Formal mathematical reading.}
Mathematical definition. This is a total EasyCrypt operation.  The exact displayed statement gives the domain, codomain, and defining equation when present.

\noindent\textbf{Role in the proof.}
Use in this chapter. It is part of the exact formal vocabulary used by subsequent lemmas and games.

\end{formaldefinition}

\begin{formaldefinition}[EasyCrypt line 75]
\noindent\textbf{EasyCrypt name.}
\begin{lstlisting}[language=EasyCrypt,basicstyle=\ttfamily\scriptsize]
keccak_lane_index
\end{lstlisting}
\noindent\textbf{Checked EasyCrypt statement.}
\begin{lstlisting}[language=EasyCrypt,basicstyle=\ttfamily\scriptsize]
op keccak_lane_index (x y : int) : int =
  (x %% 5) + 5 * (y %% 5).
\end{lstlisting}
\noindent\textbf{Formal mathematical reading.}
Mathematical definition. This is a total EasyCrypt operation.  The exact displayed statement gives the domain, codomain, and defining equation when present.

\noindent\textbf{Role in the proof.}
Use in this chapter. It is part of the exact formal vocabulary used by subsequent lemmas and games.

\end{formaldefinition}

\begin{formaldefinition}[EasyCrypt line 77]
\noindent\textbf{EasyCrypt name.}
\begin{lstlisting}[language=EasyCrypt,basicstyle=\ttfamily\scriptsize]
keccak_state_lane
\end{lstlisting}
\noindent\textbf{Checked EasyCrypt statement.}
\begin{lstlisting}[language=EasyCrypt,basicstyle=\ttfamily\scriptsize]
op keccak_state_lane (st : keccak_state) (x y : int) : keccak_lane =
  nth keccak_lane_zero st (keccak_lane_index x y).
\end{lstlisting}
\noindent\textbf{Formal mathematical reading.}
Mathematical definition. This is a total EasyCrypt operation.  The exact displayed statement gives the domain, codomain, and defining equation when present.

\noindent\textbf{Role in the proof.}
Use in this chapter. It is part of the exact formal vocabulary used by subsequent lemmas and games.

\end{formaldefinition}

\begin{formaldefinition}[EasyCrypt line 80]
\noindent\textbf{EasyCrypt name.}
\begin{lstlisting}[language=EasyCrypt,basicstyle=\ttfamily\scriptsize]
keccak_rho_offsets
\end{lstlisting}
\noindent\textbf{Checked EasyCrypt statement.}
\begin{lstlisting}[language=EasyCrypt,basicstyle=\ttfamily\scriptsize]
op keccak_rho_offsets : int list =
  0 :: 1 :: 62 :: 28 :: 27 ::
  36 :: 44 :: 6 :: 55 :: 20 ::
  3 :: 10 :: 43 :: 25 :: 39 ::
  41 :: 45 :: 15 :: 21 :: 8 ::
  18 :: 2 :: 61 :: 56 :: 14 :: [].
\end{lstlisting}
\noindent\textbf{Formal mathematical reading.}
Mathematical definition. This is a total EasyCrypt operation.  The exact displayed statement gives the domain, codomain, and defining equation when present.

\noindent\textbf{Role in the proof.}
Use in this chapter. It is part of the exact formal vocabulary used by subsequent lemmas and games.

\end{formaldefinition}

\begin{formaldefinition}[EasyCrypt line 87]
\noindent\textbf{EasyCrypt name.}
\begin{lstlisting}[language=EasyCrypt,basicstyle=\ttfamily\scriptsize]
keccak_round_constants
\end{lstlisting}
\noindent\textbf{Checked EasyCrypt statement.}
\begin{lstlisting}[language=EasyCrypt,basicstyle=\ttfamily\scriptsize]
op keccak_round_constants : int list =
  1 ::
  32898 ::
  9223372036854808714 ::
  9223372039002292224 ::
  32907 ::
  2147483649 ::
  9223372039002292353 ::
  9223372036854808585 ::
  138 ::
  136 ::
  2147516425 ::
  2147483658 ::
  2147516555 ::
  9223372036854775947 ::
  9223372036854808713 ::
  9223372036854808579 ::
  9223372036854808578 ::
  9223372036854775936 ::
  32778 ::
  9223372039002259466 ::
  9223372039002292353 ::
  9223372036854808704 ::
  2147483649 ::
  9223372039002292232 :: [].
\end{lstlisting}
\noindent\textbf{Formal mathematical reading.}
Mathematical definition. This is a total EasyCrypt operation.  The exact displayed statement gives the domain, codomain, and defining equation when present.

\noindent\textbf{Role in the proof.}
Use in this chapter. It is part of the exact formal vocabulary used by subsequent lemmas and games.

\end{formaldefinition}

\begin{formaldefinition}[EasyCrypt line 113]
\noindent\textbf{EasyCrypt name.}
\begin{lstlisting}[language=EasyCrypt,basicstyle=\ttfamily\scriptsize]
keccak_theta_c
\end{lstlisting}
\noindent\textbf{Checked EasyCrypt statement.}
\begin{lstlisting}[language=EasyCrypt,basicstyle=\ttfamily\scriptsize]
op keccak_theta_c (st : keccak_state) (x : int) : keccak_lane =
  keccak_lane_xor
    (keccak_lane_xor
      (keccak_lane_xor
        (keccak_lane_xor
          (keccak_state_lane st x 0)
          (keccak_state_lane st x 1))
        (keccak_state_lane st x 2))
      (keccak_state_lane st x 3))
    (keccak_state_lane st x 4).
\end{lstlisting}
\noindent\textbf{Formal mathematical reading.}
Mathematical definition. This is a total EasyCrypt operation.  The exact displayed statement gives the domain, codomain, and defining equation when present.

\noindent\textbf{Role in the proof.}
Use in this chapter. It is part of the exact formal vocabulary used by subsequent lemmas and games.

\end{formaldefinition}

\begin{formaldefinition}[EasyCrypt line 124]
\noindent\textbf{EasyCrypt name.}
\begin{lstlisting}[language=EasyCrypt,basicstyle=\ttfamily\scriptsize]
keccak_theta_d
\end{lstlisting}
\noindent\textbf{Checked EasyCrypt statement.}
\begin{lstlisting}[language=EasyCrypt,basicstyle=\ttfamily\scriptsize]
op keccak_theta_d (st : keccak_state) (x : int) : keccak_lane =
  keccak_lane_xor
    (keccak_theta_c st (x + 4))
    (keccak_lane_rotl (keccak_theta_c st (x + 1)) 1).
\end{lstlisting}
\noindent\textbf{Formal mathematical reading.}
Mathematical definition. This is a total EasyCrypt operation.  The exact displayed statement gives the domain, codomain, and defining equation when present.

\noindent\textbf{Role in the proof.}
Use in this chapter. It is part of the exact formal vocabulary used by subsequent lemmas and games.

\end{formaldefinition}

\begin{formaldefinition}[EasyCrypt line 129]
\noindent\textbf{EasyCrypt name.}
\begin{lstlisting}[language=EasyCrypt,basicstyle=\ttfamily\scriptsize]
keccak_theta
\end{lstlisting}
\noindent\textbf{Checked EasyCrypt statement.}
\begin{lstlisting}[language=EasyCrypt,basicstyle=\ttfamily\scriptsize]
op keccak_theta (st : keccak_state) : keccak_state =
  mkseq
    (fun i =>
       let x = i %% 5 in
       let y = i %/ 5 in
       keccak_lane_xor
         (keccak_state_lane st x y)
         (keccak_theta_d st x))
    keccak_state_lanes.
\end{lstlisting}
\noindent\textbf{Formal mathematical reading.}
Mathematical definition. This is a total EasyCrypt operation.  The exact displayed statement gives the domain, codomain, and defining equation when present.

\noindent\textbf{Role in the proof.}
Use in this chapter. It is part of the exact formal vocabulary used by subsequent lemmas and games.

\end{formaldefinition}

\begin{formaldefinition}[EasyCrypt line 139]
\noindent\textbf{EasyCrypt name.}
\begin{lstlisting}[language=EasyCrypt,basicstyle=\ttfamily\scriptsize]
keccak_rho_pi
\end{lstlisting}
\noindent\textbf{Checked EasyCrypt statement.}
\begin{lstlisting}[language=EasyCrypt,basicstyle=\ttfamily\scriptsize]
op keccak_rho_pi (st : keccak_state) : keccak_state =
  mkseq
    (fun i =>
       let x = i %% 5 in
       let y = i %/ 5 in
       let sx = (x + 3 * y) %% 5 in
       let sy = x in
       keccak_lane_rotl
         (keccak_state_lane st sx sy)
         (nth 0 keccak_rho_offsets (keccak_lane_index sx sy)))
    keccak_state_lanes.
\end{lstlisting}
\noindent\textbf{Formal mathematical reading.}
Mathematical definition. This is a total EasyCrypt operation.  The exact displayed statement gives the domain, codomain, and defining equation when present.

\noindent\textbf{Role in the proof.}
Use in this chapter. It is part of the exact formal vocabulary used by subsequent lemmas and games.

\end{formaldefinition}

\begin{formaldefinition}[EasyCrypt line 151]
\noindent\textbf{EasyCrypt name.}
\begin{lstlisting}[language=EasyCrypt,basicstyle=\ttfamily\scriptsize]
keccak_chi
\end{lstlisting}
\noindent\textbf{Checked EasyCrypt statement.}
\begin{lstlisting}[language=EasyCrypt,basicstyle=\ttfamily\scriptsize]
op keccak_chi (st : keccak_state) : keccak_state =
  mkseq
    (fun i =>
       let x = i %% 5 in
       let y = i %/ 5 in
       keccak_lane_xor
         (keccak_state_lane st x y)
         (keccak_lane_and
           (keccak_lane_not (keccak_state_lane st (x + 1) y))
           (keccak_state_lane st (x + 2) y)))
    keccak_state_lanes.
\end{lstlisting}
\noindent\textbf{Formal mathematical reading.}
Mathematical definition. This is a total EasyCrypt operation.  The exact displayed statement gives the domain, codomain, and defining equation when present.

\noindent\textbf{Role in the proof.}
Use in this chapter. It is part of the exact formal vocabulary used by subsequent lemmas and games.

\end{formaldefinition}

\begin{formaldefinition}[EasyCrypt line 163]
\noindent\textbf{EasyCrypt name.}
\begin{lstlisting}[language=EasyCrypt,basicstyle=\ttfamily\scriptsize]
keccak_iota
\end{lstlisting}
\noindent\textbf{Checked EasyCrypt statement.}
\begin{lstlisting}[language=EasyCrypt,basicstyle=\ttfamily\scriptsize]
op keccak_iota (st : keccak_state) (rc : keccak_lane) : keccak_state =
  mkseq
    (fun i =>
       if i = 0
       then keccak_lane_xor (nth keccak_lane_zero st 0) rc
       else nth keccak_lane_zero st i)
    keccak_state_lanes.
\end{lstlisting}
\noindent\textbf{Formal mathematical reading.}
Mathematical definition. This is a total EasyCrypt operation.  The exact displayed statement gives the domain, codomain, and defining equation when present.

\noindent\textbf{Role in the proof.}
Use in this chapter. It is part of the exact formal vocabulary used by subsequent lemmas and games.

\end{formaldefinition}

\begin{formaldefinition}[EasyCrypt line 171]
\noindent\textbf{EasyCrypt name.}
\begin{lstlisting}[language=EasyCrypt,basicstyle=\ttfamily\scriptsize]
keccak_round
\end{lstlisting}
\noindent\textbf{Checked EasyCrypt statement.}
\begin{lstlisting}[language=EasyCrypt,basicstyle=\ttfamily\scriptsize]
op keccak_round (st : keccak_state) (rc : keccak_lane) : keccak_state =
  keccak_iota (keccak_chi (keccak_rho_pi (keccak_theta st))) rc.
\end{lstlisting}
\noindent\textbf{Formal mathematical reading.}
Mathematical definition. This is a total EasyCrypt operation.  The exact displayed statement gives the domain, codomain, and defining equation when present.

\noindent\textbf{Role in the proof.}
Use in this chapter. It is part of the exact formal vocabulary used by subsequent lemmas and games.

\end{formaldefinition}

\begin{formaldefinition}[EasyCrypt line 174]
\noindent\textbf{EasyCrypt name.}
\begin{lstlisting}[language=EasyCrypt,basicstyle=\ttfamily\scriptsize]
keccak_f1600_lanes
\end{lstlisting}
\noindent\textbf{Checked EasyCrypt statement.}
\begin{lstlisting}[language=EasyCrypt,basicstyle=\ttfamily\scriptsize]
op keccak_f1600_lanes (st : keccak_state) : keccak_state =
  foldl
    (fun acc rc => keccak_round acc (keccak_lane_of_int rc))
    st
    keccak_round_constants.
\end{lstlisting}
\noindent\textbf{Formal mathematical reading.}
Mathematical definition. This is a total EasyCrypt operation.  The exact displayed statement gives the domain, codomain, and defining equation when present.

\noindent\textbf{Role in the proof.}
Use in this chapter. It is part of the exact formal vocabulary used by subsequent lemmas and games.

\end{formaldefinition}

\begin{formaldefinition}[EasyCrypt line 180]
\noindent\textbf{EasyCrypt name.}
\begin{lstlisting}[language=EasyCrypt,basicstyle=\ttfamily\scriptsize]
keccak_f1600_bytes
\end{lstlisting}
\noindent\textbf{Checked EasyCrypt statement.}
\begin{lstlisting}[language=EasyCrypt,basicstyle=\ttfamily\scriptsize]
op keccak_f1600_bytes (st : keccak_byte list) : keccak_byte list =
  keccak_bytes_of_lanes (keccak_f1600_lanes (keccak_lanes_of_bytes st)).
\end{lstlisting}
\noindent\textbf{Formal mathematical reading.}
Mathematical definition. This is a total EasyCrypt operation.  The exact displayed statement gives the domain, codomain, and defining equation when present.

\noindent\textbf{Role in the proof.}
Use in this chapter. It is part of the exact formal vocabulary used by subsequent lemmas and games.

\end{formaldefinition}

\begin{formaldefinition}[EasyCrypt line 381]
\noindent\textbf{EasyCrypt name.}
\begin{lstlisting}[language=EasyCrypt,basicstyle=\ttfamily\scriptsize]
keccak_theta_c_bit_index_value
\end{lstlisting}
\noindent\textbf{Checked EasyCrypt statement.}
\begin{lstlisting}[language=EasyCrypt,basicstyle=\ttfamily\scriptsize]
op keccak_theta_c_bit_index_value
   (st : keccak_state) (x bit : int) : bool =
  ((((keccak_lane_bit
        (nth keccak_lane_zero st (keccak_lane_index x 0)) bit <>
      keccak_lane_bit
        (nth keccak_lane_zero st (keccak_lane_index x 1)) bit) <>
     keccak_lane_bit
       (nth keccak_lane_zero st (keccak_lane_index x 2)) bit) <>
    keccak_lane_bit
      (nth keccak_lane_zero st (keccak_lane_index x 3)) bit) <>
   keccak_lane_bit
     (nth keccak_lane_zero st (keccak_lane_index x 4)) bit).
\end{lstlisting}
\noindent\textbf{Formal mathematical reading.}
Mathematical definition. This is a Boolean predicate.  The exact displayed EasyCrypt statement gives its arguments; mathematically it is an event, invariant, support condition, or well-formedness predicate over those arguments.

\noindent\textbf{Role in the proof.}
Use in this chapter. It is part of the exact formal vocabulary used by subsequent lemmas and games.

\end{formaldefinition}

\begin{formaldefinition}[EasyCrypt line 394]
\noindent\textbf{EasyCrypt name.}
\begin{lstlisting}[language=EasyCrypt,basicstyle=\ttfamily\scriptsize]
keccak_theta_d_bit_index_value
\end{lstlisting}
\noindent\textbf{Checked EasyCrypt statement.}
\begin{lstlisting}[language=EasyCrypt,basicstyle=\ttfamily\scriptsize]
op keccak_theta_d_bit_index_value
   (st : keccak_state) (x bit : int) : bool =
  keccak_theta_c_bit_index_value st (x + 4) bit <>
  keccak_theta_c_bit_index_value
    st
    (x + 1)
    ((bit - 1) %% keccak_lane_bits).
\end{lstlisting}
\noindent\textbf{Formal mathematical reading.}
Mathematical definition. This is a Boolean predicate.  The exact displayed EasyCrypt statement gives its arguments; mathematically it is an event, invariant, support condition, or well-formedness predicate over those arguments.

\noindent\textbf{Role in the proof.}
Use in this chapter. It is part of the exact formal vocabulary used by subsequent lemmas and games.

\end{formaldefinition}

\subsubsection*{Lemmas, Theorems, Relational Equivalences, and Their Proof Dependencies}
\begin{formaltheorem}[EasyCrypt line 183]
\noindent\textbf{EasyCrypt name.}
\begin{lstlisting}[language=EasyCrypt,basicstyle=\ttfamily\scriptsize]
keccak_lane_bitsE
\end{lstlisting}
\noindent\textbf{Checked EasyCrypt statement.}
\begin{lstlisting}[language=EasyCrypt,basicstyle=\ttfamily\scriptsize]
lemma keccak_lane_bitsE :
  keccak_lane_bits = 64.
\end{lstlisting}
\noindent\textbf{Formal mathematical reading.}
Mathematical theorem. This is an equality or definitional expansion used for rewriting and normalization.

\noindent\textbf{Role in the proof.}
Use in this chapter. It is a reusable checked theorem cited by later proof scripts.

\noindent\textbf{Proof-script interpretation.}
No proof block was collected for this item.

\end{formaltheorem}

\begin{formaltheorem}[EasyCrypt line 187]
\noindent\textbf{EasyCrypt name.}
\begin{lstlisting}[language=EasyCrypt,basicstyle=\ttfamily\scriptsize]
keccak_state_lanesE
\end{lstlisting}
\noindent\textbf{Checked EasyCrypt statement.}
\begin{lstlisting}[language=EasyCrypt,basicstyle=\ttfamily\scriptsize]
lemma keccak_state_lanesE :
  keccak_state_lanes = 25.
\end{lstlisting}
\noindent\textbf{Formal mathematical reading.}
Mathematical theorem. This is an equality or definitional expansion used for rewriting and normalization.

\noindent\textbf{Role in the proof.}
Use in this chapter. It preserves an invariant across an oracle call, adversary call, or game procedure.

\noindent\textbf{Proof-script interpretation.}
No proof block was collected for this item.

\end{formaltheorem}

\begin{formaltheorem}[EasyCrypt line 191]
\noindent\textbf{EasyCrypt name.}
\begin{lstlisting}[language=EasyCrypt,basicstyle=\ttfamily\scriptsize]
keccak_state_bytesE
\end{lstlisting}
\noindent\textbf{Checked EasyCrypt statement.}
\begin{lstlisting}[language=EasyCrypt,basicstyle=\ttfamily\scriptsize]
lemma keccak_state_bytesE :
  keccak_state_bytes = 200.
\end{lstlisting}
\noindent\textbf{Formal mathematical reading.}
Mathematical theorem. This is an equality or definitional expansion used for rewriting and normalization.

\noindent\textbf{Role in the proof.}
Use in this chapter. It preserves an invariant across an oracle call, adversary call, or game procedure.

\noindent\textbf{Proof-script interpretation.}
No proof block was collected for this item.

\end{formaltheorem}

\begin{formaltheorem}[EasyCrypt line 195]
\noindent\textbf{EasyCrypt name.}
\begin{lstlisting}[language=EasyCrypt,basicstyle=\ttfamily\scriptsize]
keccak_lane_zero_wf
\end{lstlisting}
\noindent\textbf{Checked EasyCrypt statement.}
\begin{lstlisting}[language=EasyCrypt,basicstyle=\ttfamily\scriptsize]
lemma keccak_lane_zero_wf :
  keccak_lane_wf keccak_lane_zero.
\end{lstlisting}
\noindent\textbf{Formal mathematical reading.}
Mathematical theorem. This lemma is a universally quantified proposition over the variables shown in the EasyCrypt statement.

\noindent\textbf{Role in the proof.}
Use in this chapter. It proves a structural correctness fact needed by later extraction or verification arguments.

\noindent\textbf{Proof-script interpretation.}
No proof block was collected for this item.

\end{formaltheorem}

\begin{formaltheorem}[EasyCrypt line 199]
\noindent\textbf{EasyCrypt name.}
\begin{lstlisting}[language=EasyCrypt,basicstyle=\ttfamily\scriptsize]
keccak_lane_xor_wf
\end{lstlisting}
\noindent\textbf{Checked EasyCrypt statement.}
\begin{lstlisting}[language=EasyCrypt,basicstyle=\ttfamily\scriptsize]
lemma keccak_lane_xor_wf a b :
  keccak_lane_wf (keccak_lane_xor a b).
\end{lstlisting}
\noindent\textbf{Formal mathematical reading.}
Mathematical theorem. This lemma is a universally quantified proposition over the variables shown in the EasyCrypt statement.

\noindent\textbf{Role in the proof.}
Use in this chapter. It proves a structural correctness fact needed by later extraction or verification arguments.

\noindent\textbf{Proof-script interpretation.}
No proof block was collected for this item.

\end{formaltheorem}

\begin{formaltheorem}[EasyCrypt line 203]
\noindent\textbf{EasyCrypt name.}
\begin{lstlisting}[language=EasyCrypt,basicstyle=\ttfamily\scriptsize]
keccak_lane_and_wf
\end{lstlisting}
\noindent\textbf{Checked EasyCrypt statement.}
\begin{lstlisting}[language=EasyCrypt,basicstyle=\ttfamily\scriptsize]
lemma keccak_lane_and_wf a b :
  keccak_lane_wf (keccak_lane_and a b).
\end{lstlisting}
\noindent\textbf{Formal mathematical reading.}
Mathematical theorem. This lemma is a universally quantified proposition over the variables shown in the EasyCrypt statement.

\noindent\textbf{Role in the proof.}
Use in this chapter. It proves a structural correctness fact needed by later extraction or verification arguments.

\noindent\textbf{Proof-script interpretation.}
No proof block was collected for this item.

\end{formaltheorem}

\begin{formaltheorem}[EasyCrypt line 207]
\noindent\textbf{EasyCrypt name.}
\begin{lstlisting}[language=EasyCrypt,basicstyle=\ttfamily\scriptsize]
keccak_lane_not_wf
\end{lstlisting}
\noindent\textbf{Checked EasyCrypt statement.}
\begin{lstlisting}[language=EasyCrypt,basicstyle=\ttfamily\scriptsize]
lemma keccak_lane_not_wf a :
  keccak_lane_wf (keccak_lane_not a).
\end{lstlisting}
\noindent\textbf{Formal mathematical reading.}
Mathematical theorem. This lemma is a universally quantified proposition over the variables shown in the EasyCrypt statement.

\noindent\textbf{Role in the proof.}
Use in this chapter. It proves a structural correctness fact needed by later extraction or verification arguments.

\noindent\textbf{Proof-script interpretation.}
No proof block was collected for this item.

\end{formaltheorem}

\begin{formaltheorem}[EasyCrypt line 211]
\noindent\textbf{EasyCrypt name.}
\begin{lstlisting}[language=EasyCrypt,basicstyle=\ttfamily\scriptsize]
keccak_lane_rotl_wf
\end{lstlisting}
\noindent\textbf{Checked EasyCrypt statement.}
\begin{lstlisting}[language=EasyCrypt,basicstyle=\ttfamily\scriptsize]
lemma keccak_lane_rotl_wf a offset :
  keccak_lane_wf (keccak_lane_rotl a offset).
\end{lstlisting}
\noindent\textbf{Formal mathematical reading.}
Mathematical theorem. This lemma is a universally quantified proposition over the variables shown in the EasyCrypt statement.

\noindent\textbf{Role in the proof.}
Use in this chapter. It proves a structural correctness fact needed by later extraction or verification arguments.

\noindent\textbf{Proof-script interpretation.}
No proof block was collected for this item.

\end{formaltheorem}

\begin{formaltheorem}[EasyCrypt line 215]
\noindent\textbf{EasyCrypt name.}
\begin{lstlisting}[language=EasyCrypt,basicstyle=\ttfamily\scriptsize]
keccak_lane_of_int_wf
\end{lstlisting}
\noindent\textbf{Checked EasyCrypt statement.}
\begin{lstlisting}[language=EasyCrypt,basicstyle=\ttfamily\scriptsize]
lemma keccak_lane_of_int_wf x :
  keccak_lane_wf (keccak_lane_of_int x).
\end{lstlisting}
\noindent\textbf{Formal mathematical reading.}
Mathematical theorem. This lemma is a universally quantified proposition over the variables shown in the EasyCrypt statement.

\noindent\textbf{Role in the proof.}
Use in this chapter. It proves a structural correctness fact needed by later extraction or verification arguments.

\noindent\textbf{Proof-script interpretation.}
No proof block was collected for this item.

\end{formaltheorem}

\begin{formaltheorem}[EasyCrypt line 219]
\noindent\textbf{EasyCrypt name.}
\begin{lstlisting}[language=EasyCrypt,basicstyle=\ttfamily\scriptsize]
keccak_lane_mod_range
\end{lstlisting}
\noindent\textbf{Checked EasyCrypt statement.}
\begin{lstlisting}[language=EasyCrypt,basicstyle=\ttfamily\scriptsize]
lemma keccak_lane_mod_range x :
  0 <= x %% keccak_lane_bits < keccak_lane_bits.
\end{lstlisting}
\noindent\textbf{Formal mathematical reading.}
Mathematical theorem. This is an inequality, usually an arithmetic loss bound, event inclusion bound, or monotonicity fact used to compose probabilities.

\noindent\textbf{Role in the proof.}
Use in this chapter. It is a reusable checked theorem cited by later proof scripts.

\noindent\textbf{Proof-script interpretation.}
No proof block was collected for this item.

\end{formaltheorem}

\begin{formaltheorem}[EasyCrypt line 223]
\noindent\textbf{EasyCrypt name.}
\begin{lstlisting}[language=EasyCrypt,basicstyle=\ttfamily\scriptsize]
keccak_word_modulus_pos
\end{lstlisting}
\noindent\textbf{Checked EasyCrypt statement.}
\begin{lstlisting}[language=EasyCrypt,basicstyle=\ttfamily\scriptsize]
lemma keccak_word_modulus_pos :
  0 < 2 ^ keccak_lane_bits.
\end{lstlisting}
\noindent\textbf{Formal mathematical reading.}
Mathematical theorem. This lemma is a universally quantified proposition over the variables shown in the EasyCrypt statement.

\noindent\textbf{Role in the proof.}
Use in this chapter. It is a reusable checked theorem cited by later proof scripts.

\noindent\textbf{Proof-script interpretation.}
No proof block was collected for this item.

\end{formaltheorem}

\begin{formaltheorem}[EasyCrypt line 227]
\noindent\textbf{EasyCrypt name.}
\begin{lstlisting}[language=EasyCrypt,basicstyle=\ttfamily\scriptsize]
keccak_word_norm_small
\end{lstlisting}
\noindent\textbf{Checked EasyCrypt statement.}
\begin{lstlisting}[language=EasyCrypt,basicstyle=\ttfamily\scriptsize]
lemma keccak_word_norm_small w :
  0 <= w < 2 ^ keccak_lane_bits =>
  keccak_word_norm w = w.
\end{lstlisting}
\noindent\textbf{Formal mathematical reading.}
Mathematical theorem. This is an inequality, usually an arithmetic loss bound, event inclusion bound, or monotonicity fact used to compose probabilities.

\noindent\textbf{Role in the proof.}
Use in this chapter. It is a reusable checked theorem cited by later proof scripts.

\noindent\textbf{Proof-script interpretation.}
Proof-script reading. The machine proof decomposes the theorem into rewrites, program-logic obligations, relational equivalences, and arithmetic side conditions.
\emph{Main tactics visible in the proof script:} \ec{rewrite}.
\emph{Named identifiers syntactically cited in the proof block:}
\begin{lstlisting}[language=EasyCrypt,basicstyle=\ttfamily\scriptsize]
keccak_word_norm
pmod_small
w_range
\end{lstlisting}

\end{formaltheorem}

\begin{formaltheorem}[EasyCrypt line 235]
\noindent\textbf{EasyCrypt name.}
\begin{lstlisting}[language=EasyCrypt,basicstyle=\ttfamily\scriptsize]
keccak_pow2_0E
\end{lstlisting}
\noindent\textbf{Checked EasyCrypt statement.}
\begin{lstlisting}[language=EasyCrypt,basicstyle=\ttfamily\scriptsize]
lemma keccak_pow2_0E :
  2 ^ 0 = 1.
\end{lstlisting}
\noindent\textbf{Formal mathematical reading.}
Mathematical theorem. This is an equality or definitional expansion used for rewriting and normalization.

\noindent\textbf{Role in the proof.}
Use in this chapter. It is a reusable checked theorem cited by later proof scripts.

\noindent\textbf{Proof-script interpretation.}
No proof block was collected for this item.

\end{formaltheorem}

\begin{formaltheorem}[EasyCrypt line 239]
\noindent\textbf{EasyCrypt name.}
\begin{lstlisting}[language=EasyCrypt,basicstyle=\ttfamily\scriptsize]
keccak_pow2_1E
\end{lstlisting}
\noindent\textbf{Checked EasyCrypt statement.}
\begin{lstlisting}[language=EasyCrypt,basicstyle=\ttfamily\scriptsize]
lemma keccak_pow2_1E :
  2 ^ 1 = 2.
\end{lstlisting}
\noindent\textbf{Formal mathematical reading.}
Mathematical theorem. This is an equality or definitional expansion used for rewriting and normalization.

\noindent\textbf{Role in the proof.}
Use in this chapter. It is a reusable checked theorem cited by later proof scripts.

\noindent\textbf{Proof-script interpretation.}
No proof block was collected for this item.

\end{formaltheorem}

\begin{formaltheorem}[EasyCrypt line 243]
\noindent\textbf{EasyCrypt name.}
\begin{lstlisting}[language=EasyCrypt,basicstyle=\ttfamily\scriptsize]
keccak_pow2_2E
\end{lstlisting}
\noindent\textbf{Checked EasyCrypt statement.}
\begin{lstlisting}[language=EasyCrypt,basicstyle=\ttfamily\scriptsize]
lemma keccak_pow2_2E :
  2 ^ 2 = 4.
\end{lstlisting}
\noindent\textbf{Formal mathematical reading.}
Mathematical theorem. This is an equality or definitional expansion used for rewriting and normalization.

\noindent\textbf{Role in the proof.}
Use in this chapter. It is a reusable checked theorem cited by later proof scripts.

\noindent\textbf{Proof-script interpretation.}
No proof block was collected for this item.

\end{formaltheorem}

\begin{formaltheorem}[EasyCrypt line 247]
\noindent\textbf{EasyCrypt name.}
\begin{lstlisting}[language=EasyCrypt,basicstyle=\ttfamily\scriptsize]
keccak_pow2_3E
\end{lstlisting}
\noindent\textbf{Checked EasyCrypt statement.}
\begin{lstlisting}[language=EasyCrypt,basicstyle=\ttfamily\scriptsize]
lemma keccak_pow2_3E :
  2 ^ 3 = 8.
\end{lstlisting}
\noindent\textbf{Formal mathematical reading.}
Mathematical theorem. This is an equality or definitional expansion used for rewriting and normalization.

\noindent\textbf{Role in the proof.}
Use in this chapter. It is a reusable checked theorem cited by later proof scripts.

\noindent\textbf{Proof-script interpretation.}
Proof-script reading. The machine proof decomposes the theorem into rewrites, program-logic obligations, relational equivalences, and arithmetic side conditions.
\emph{Main tactics visible in the proof script:} \ec{rewrite}.
\emph{Named identifiers syntactically cited in the proof block:}
\begin{lstlisting}[language=EasyCrypt,basicstyle=\ttfamily\scriptsize]
exprS
keccak_pow2_2E
\end{lstlisting}

\end{formaltheorem}

\begin{formaltheorem}[EasyCrypt line 254]
\noindent\textbf{EasyCrypt name.}
\begin{lstlisting}[language=EasyCrypt,basicstyle=\ttfamily\scriptsize]
keccak_pow2_4E
\end{lstlisting}
\noindent\textbf{Checked EasyCrypt statement.}
\begin{lstlisting}[language=EasyCrypt,basicstyle=\ttfamily\scriptsize]
lemma keccak_pow2_4E :
  2 ^ 4 = 16.
\end{lstlisting}
\noindent\textbf{Formal mathematical reading.}
Mathematical theorem. This is an equality or definitional expansion used for rewriting and normalization.

\noindent\textbf{Role in the proof.}
Use in this chapter. It is a reusable checked theorem cited by later proof scripts.

\noindent\textbf{Proof-script interpretation.}
Proof-script reading. The machine proof decomposes the theorem into rewrites, program-logic obligations, relational equivalences, and arithmetic side conditions.
\emph{Main tactics visible in the proof script:} \ec{rewrite}.
\emph{Named identifiers syntactically cited in the proof block:}
\begin{lstlisting}[language=EasyCrypt,basicstyle=\ttfamily\scriptsize]
exprS
keccak_pow2_3E
\end{lstlisting}

\end{formaltheorem}

\begin{formaltheorem}[EasyCrypt line 261]
\noindent\textbf{EasyCrypt name.}
\begin{lstlisting}[language=EasyCrypt,basicstyle=\ttfamily\scriptsize]
keccak_pow2_5E
\end{lstlisting}
\noindent\textbf{Checked EasyCrypt statement.}
\begin{lstlisting}[language=EasyCrypt,basicstyle=\ttfamily\scriptsize]
lemma keccak_pow2_5E :
  2 ^ 5 = 32.
\end{lstlisting}
\noindent\textbf{Formal mathematical reading.}
Mathematical theorem. This is an equality or definitional expansion used for rewriting and normalization.

\noindent\textbf{Role in the proof.}
Use in this chapter. It is a reusable checked theorem cited by later proof scripts.

\noindent\textbf{Proof-script interpretation.}
Proof-script reading. The machine proof decomposes the theorem into rewrites, program-logic obligations, relational equivalences, and arithmetic side conditions.
\emph{Main tactics visible in the proof script:} \ec{rewrite}.
\emph{Named identifiers syntactically cited in the proof block:}
\begin{lstlisting}[language=EasyCrypt,basicstyle=\ttfamily\scriptsize]
exprS
keccak_pow2_4E
\end{lstlisting}

\end{formaltheorem}

\begin{formaltheorem}[EasyCrypt line 268]
\noindent\textbf{EasyCrypt name.}
\begin{lstlisting}[language=EasyCrypt,basicstyle=\ttfamily\scriptsize]
keccak_pow2_6E
\end{lstlisting}
\noindent\textbf{Checked EasyCrypt statement.}
\begin{lstlisting}[language=EasyCrypt,basicstyle=\ttfamily\scriptsize]
lemma keccak_pow2_6E :
  2 ^ 6 = 64.
\end{lstlisting}
\noindent\textbf{Formal mathematical reading.}
Mathematical theorem. This is an equality or definitional expansion used for rewriting and normalization.

\noindent\textbf{Role in the proof.}
Use in this chapter. It is a reusable checked theorem cited by later proof scripts.

\noindent\textbf{Proof-script interpretation.}
Proof-script reading. The machine proof decomposes the theorem into rewrites, program-logic obligations, relational equivalences, and arithmetic side conditions.
\emph{Main tactics visible in the proof script:} \ec{rewrite}.
\emph{Named identifiers syntactically cited in the proof block:}
\begin{lstlisting}[language=EasyCrypt,basicstyle=\ttfamily\scriptsize]
exprS
keccak_pow2_5E
\end{lstlisting}

\end{formaltheorem}

\begin{formaltheorem}[EasyCrypt line 275]
\noindent\textbf{EasyCrypt name.}
\begin{lstlisting}[language=EasyCrypt,basicstyle=\ttfamily\scriptsize]
keccak_pow2_7E
\end{lstlisting}
\noindent\textbf{Checked EasyCrypt statement.}
\begin{lstlisting}[language=EasyCrypt,basicstyle=\ttfamily\scriptsize]
lemma keccak_pow2_7E :
  2 ^ 7 = 128.
\end{lstlisting}
\noindent\textbf{Formal mathematical reading.}
Mathematical theorem. This is an equality or definitional expansion used for rewriting and normalization.

\noindent\textbf{Role in the proof.}
Use in this chapter. It is a reusable checked theorem cited by later proof scripts.

\noindent\textbf{Proof-script interpretation.}
Proof-script reading. The machine proof decomposes the theorem into rewrites, program-logic obligations, relational equivalences, and arithmetic side conditions.
\emph{Main tactics visible in the proof script:} \ec{rewrite}.
\emph{Named identifiers syntactically cited in the proof block:}
\begin{lstlisting}[language=EasyCrypt,basicstyle=\ttfamily\scriptsize]
exprS
keccak_pow2_6E
\end{lstlisting}

\end{formaltheorem}

\begin{formaltheorem}[EasyCrypt line 282]
\noindent\textbf{EasyCrypt name.}
\begin{lstlisting}[language=EasyCrypt,basicstyle=\ttfamily\scriptsize]
keccak_bytes_to_lane_wf
\end{lstlisting}
\noindent\textbf{Checked EasyCrypt statement.}
\begin{lstlisting}[language=EasyCrypt,basicstyle=\ttfamily\scriptsize]
lemma keccak_bytes_to_lane_wf bs lane_index :
  keccak_lane_wf (keccak_bytes_to_lane bs lane_index).
\end{lstlisting}
\noindent\textbf{Formal mathematical reading.}
Mathematical theorem. This lemma is a universally quantified proposition over the variables shown in the EasyCrypt statement.

\noindent\textbf{Role in the proof.}
Use in this chapter. It proves a structural correctness fact needed by later extraction or verification arguments.

\noindent\textbf{Proof-script interpretation.}
No proof block was collected for this item.

\end{formaltheorem}

\begin{formaltheorem}[EasyCrypt line 286]
\noindent\textbf{EasyCrypt name.}
\begin{lstlisting}[language=EasyCrypt,basicstyle=\ttfamily\scriptsize]
keccak_rho_offsets_size
\end{lstlisting}
\noindent\textbf{Checked EasyCrypt statement.}
\begin{lstlisting}[language=EasyCrypt,basicstyle=\ttfamily\scriptsize]
lemma keccak_rho_offsets_size :
  size keccak_rho_offsets = keccak_state_lanes.
\end{lstlisting}
\noindent\textbf{Formal mathematical reading.}
Mathematical theorem. This is an equality or definitional expansion used for rewriting and normalization.

\noindent\textbf{Role in the proof.}
Use in this chapter. It is a reusable checked theorem cited by later proof scripts.

\noindent\textbf{Proof-script interpretation.}
No proof block was collected for this item.

\end{formaltheorem}

\begin{formaltheorem}[EasyCrypt line 290]
\noindent\textbf{EasyCrypt name.}
\begin{lstlisting}[language=EasyCrypt,basicstyle=\ttfamily\scriptsize]
keccak_round_constants_size
\end{lstlisting}
\noindent\textbf{Checked EasyCrypt statement.}
\begin{lstlisting}[language=EasyCrypt,basicstyle=\ttfamily\scriptsize]
lemma keccak_round_constants_size :
  size keccak_round_constants = 24.
\end{lstlisting}
\noindent\textbf{Formal mathematical reading.}
Mathematical theorem. This is an equality or definitional expansion used for rewriting and normalization.

\noindent\textbf{Role in the proof.}
Use in this chapter. It is a reusable checked theorem cited by later proof scripts.

\noindent\textbf{Proof-script interpretation.}
No proof block was collected for this item.

\end{formaltheorem}

\begin{formaltheorem}[EasyCrypt line 294]
\noindent\textbf{EasyCrypt name.}
\begin{lstlisting}[language=EasyCrypt,basicstyle=\ttfamily\scriptsize]
keccak_lanes_of_bytes_size
\end{lstlisting}
\noindent\textbf{Checked EasyCrypt statement.}
\begin{lstlisting}[language=EasyCrypt,basicstyle=\ttfamily\scriptsize]
lemma keccak_lanes_of_bytes_size bs :
  size (keccak_lanes_of_bytes bs) = keccak_state_lanes.
\end{lstlisting}
\noindent\textbf{Formal mathematical reading.}
Mathematical theorem. This is an equality or definitional expansion used for rewriting and normalization.

\noindent\textbf{Role in the proof.}
Use in this chapter. It is a reusable checked theorem cited by later proof scripts.

\noindent\textbf{Proof-script interpretation.}
No proof block was collected for this item.

\end{formaltheorem}

\begin{formaltheorem}[EasyCrypt line 298]
\noindent\textbf{EasyCrypt name.}
\begin{lstlisting}[language=EasyCrypt,basicstyle=\ttfamily\scriptsize]
keccak_lanes_of_bytes_lane_wf
\end{lstlisting}
\noindent\textbf{Checked EasyCrypt statement.}
\begin{lstlisting}[language=EasyCrypt,basicstyle=\ttfamily\scriptsize]
lemma keccak_lanes_of_bytes_lane_wf bs lane :
  0 <= lane < keccak_state_lanes =>
  keccak_lane_wf
    (nth keccak_lane_zero (keccak_lanes_of_bytes bs) lane).
\end{lstlisting}
\noindent\textbf{Formal mathematical reading.}
Mathematical theorem. This is an inequality, usually an arithmetic loss bound, event inclusion bound, or monotonicity fact used to compose probabilities.

\noindent\textbf{Role in the proof.}
Use in this chapter. It proves a structural correctness fact needed by later extraction or verification arguments.

\noindent\textbf{Proof-script interpretation.}
Proof-script reading. The machine proof decomposes the theorem into rewrites, program-logic obligations, relational equivalences, and arithmetic side conditions.
\emph{Main tactics visible in the proof script:} \ec{rewrite}, \ec{smt}, \ec{apply}.
\emph{Named identifiers syntactically cited in the proof block:}
\begin{lstlisting}[language=EasyCrypt,basicstyle=\ttfamily\scriptsize]
keccak_bytes_to_lane_wf
keccak_lanes_of_bytes
lane_range
nth_mkseq
\end{lstlisting}

\end{formaltheorem}

\begin{formaltheorem}[EasyCrypt line 308]
\noindent\textbf{EasyCrypt name.}
\begin{lstlisting}[language=EasyCrypt,basicstyle=\ttfamily\scriptsize]
keccak_lanes_of_bytes_bitE
\end{lstlisting}
\noindent\textbf{Checked EasyCrypt statement.}
\begin{lstlisting}[language=EasyCrypt,basicstyle=\ttfamily\scriptsize]
lemma keccak_lanes_of_bytes_bitE bs lane bit :
  0 <= lane < keccak_state_lanes =>
  0 <= bit < keccak_lane_bits =>
  keccak_lane_bit
    (nth keccak_lane_zero (keccak_lanes_of_bytes bs) lane)
    bit =
  keccak_byte_bit
    (nth 0 bs (8 * lane + bit %/ 8))
    (bit %% 8).
\end{lstlisting}
\noindent\textbf{Formal mathematical reading.}
Mathematical theorem. This is an inequality, usually an arithmetic loss bound, event inclusion bound, or monotonicity fact used to compose probabilities.

\noindent\textbf{Role in the proof.}
Use in this chapter. It is a reusable checked theorem cited by later proof scripts.

\noindent\textbf{Proof-script interpretation.}
Proof-script reading. The machine proof decomposes the theorem into rewrites, program-logic obligations, relational equivalences, and arithmetic side conditions.
\emph{Main tactics visible in the proof script:} \ec{rewrite}, \ec{smt}.
\emph{Named identifiers syntactically cited in the proof block:}
\begin{lstlisting}[language=EasyCrypt,basicstyle=\ttfamily\scriptsize]
bit_range
keccak_bytes_to_lane
keccak_lane_bit
keccak_lanes_of_bytes
lane_range
nth_mkseq
\end{lstlisting}

\end{formaltheorem}

\begin{formaltheorem}[EasyCrypt line 324]
\noindent\textbf{EasyCrypt name.}
\begin{lstlisting}[language=EasyCrypt,basicstyle=\ttfamily\scriptsize]
keccak_lane_xor_bitE
\end{lstlisting}
\noindent\textbf{Checked EasyCrypt statement.}
\begin{lstlisting}[language=EasyCrypt,basicstyle=\ttfamily\scriptsize]
lemma keccak_lane_xor_bitE a b bit :
  0 <= bit < keccak_lane_bits =>
  keccak_lane_bit (keccak_lane_xor a b) bit =
  (keccak_lane_bit a bit <> keccak_lane_bit b bit).
\end{lstlisting}
\noindent\textbf{Formal mathematical reading.}
Mathematical theorem. This is an inequality, usually an arithmetic loss bound, event inclusion bound, or monotonicity fact used to compose probabilities.

\noindent\textbf{Role in the proof.}
Use in this chapter. It is a reusable checked theorem cited by later proof scripts.

\noindent\textbf{Proof-script interpretation.}
Proof-script reading. The machine proof decomposes the theorem into rewrites, program-logic obligations, relational equivalences, and arithmetic side conditions.
\emph{Main tactics visible in the proof script:} \ec{rewrite}, \ec{smt}.
\emph{Named identifiers syntactically cited in the proof block:}
\begin{lstlisting}[language=EasyCrypt,basicstyle=\ttfamily\scriptsize]
bit_range
keccak_lane_bit
keccak_lane_xor
nth_mkseq
\end{lstlisting}

\end{formaltheorem}

\begin{formaltheorem}[EasyCrypt line 334]
\noindent\textbf{EasyCrypt name.}
\begin{lstlisting}[language=EasyCrypt,basicstyle=\ttfamily\scriptsize]
keccak_lane_and_bitE
\end{lstlisting}
\noindent\textbf{Checked EasyCrypt statement.}
\begin{lstlisting}[language=EasyCrypt,basicstyle=\ttfamily\scriptsize]
lemma keccak_lane_and_bitE a b bit :
  0 <= bit < keccak_lane_bits =>
  keccak_lane_bit (keccak_lane_and a b) bit =
  (keccak_lane_bit a bit /\ keccak_lane_bit b bit).
\end{lstlisting}
\noindent\textbf{Formal mathematical reading.}
Mathematical theorem. This is an inequality, usually an arithmetic loss bound, event inclusion bound, or monotonicity fact used to compose probabilities.

\noindent\textbf{Role in the proof.}
Use in this chapter. It is a reusable checked theorem cited by later proof scripts.

\noindent\textbf{Proof-script interpretation.}
Proof-script reading. The machine proof decomposes the theorem into rewrites, program-logic obligations, relational equivalences, and arithmetic side conditions.
\emph{Main tactics visible in the proof script:} \ec{rewrite}, \ec{smt}.
\emph{Named identifiers syntactically cited in the proof block:}
\begin{lstlisting}[language=EasyCrypt,basicstyle=\ttfamily\scriptsize]
bit_range
keccak_lane_and
keccak_lane_bit
nth_mkseq
\end{lstlisting}

\end{formaltheorem}

\begin{formaltheorem}[EasyCrypt line 344]
\noindent\textbf{EasyCrypt name.}
\begin{lstlisting}[language=EasyCrypt,basicstyle=\ttfamily\scriptsize]
keccak_lane_not_bitE
\end{lstlisting}
\noindent\textbf{Checked EasyCrypt statement.}
\begin{lstlisting}[language=EasyCrypt,basicstyle=\ttfamily\scriptsize]
lemma keccak_lane_not_bitE a bit :
  0 <= bit < keccak_lane_bits =>
  keccak_lane_bit (keccak_lane_not a) bit =
  (if keccak_lane_bit a bit then false else true).
\end{lstlisting}
\noindent\textbf{Formal mathematical reading.}
Mathematical theorem. This is an inequality, usually an arithmetic loss bound, event inclusion bound, or monotonicity fact used to compose probabilities.

\noindent\textbf{Role in the proof.}
Use in this chapter. It is a reusable checked theorem cited by later proof scripts.

\noindent\textbf{Proof-script interpretation.}
Proof-script reading. The machine proof decomposes the theorem into rewrites, program-logic obligations, relational equivalences, and arithmetic side conditions.
\emph{Main tactics visible in the proof script:} \ec{rewrite}, \ec{smt}.
\emph{Named identifiers syntactically cited in the proof block:}
\begin{lstlisting}[language=EasyCrypt,basicstyle=\ttfamily\scriptsize]
bit_range
keccak_lane_bit
keccak_lane_not
nth_mkseq
\end{lstlisting}

\end{formaltheorem}

\begin{formaltheorem}[EasyCrypt line 354]
\noindent\textbf{EasyCrypt name.}
\begin{lstlisting}[language=EasyCrypt,basicstyle=\ttfamily\scriptsize]
keccak_lane_rotl_bitE
\end{lstlisting}
\noindent\textbf{Checked EasyCrypt statement.}
\begin{lstlisting}[language=EasyCrypt,basicstyle=\ttfamily\scriptsize]
lemma keccak_lane_rotl_bitE a offset bit :
  0 <= bit < keccak_lane_bits =>
  keccak_lane_bit (keccak_lane_rotl a offset) bit =
  keccak_lane_bit a ((bit - offset) %% keccak_lane_bits).
\end{lstlisting}
\noindent\textbf{Formal mathematical reading.}
Mathematical theorem. This is an inequality, usually an arithmetic loss bound, event inclusion bound, or monotonicity fact used to compose probabilities.

\noindent\textbf{Role in the proof.}
Use in this chapter. It is a reusable checked theorem cited by later proof scripts.

\noindent\textbf{Proof-script interpretation.}
Proof-script reading. The machine proof decomposes the theorem into rewrites, program-logic obligations, relational equivalences, and arithmetic side conditions.
\emph{Main tactics visible in the proof script:} \ec{rewrite}, \ec{smt}.
\emph{Named identifiers syntactically cited in the proof block:}
\begin{lstlisting}[language=EasyCrypt,basicstyle=\ttfamily\scriptsize]
bit_range
keccak_lane_bit
keccak_lane_rotl
nth_mkseq
\end{lstlisting}

\end{formaltheorem}

\begin{formaltheorem}[EasyCrypt line 364]
\noindent\textbf{EasyCrypt name.}
\begin{lstlisting}[language=EasyCrypt,basicstyle=\ttfamily\scriptsize]
keccak_lane_of_int_bitE
\end{lstlisting}
\noindent\textbf{Checked EasyCrypt statement.}
\begin{lstlisting}[language=EasyCrypt,basicstyle=\ttfamily\scriptsize]
lemma keccak_lane_of_int_bitE x bit :
  0 <= bit < keccak_lane_bits =>
  keccak_lane_bit (keccak_lane_of_int x) bit =
  keccak_word_bit x bit.
\end{lstlisting}
\noindent\textbf{Formal mathematical reading.}
Mathematical theorem. This is an inequality, usually an arithmetic loss bound, event inclusion bound, or monotonicity fact used to compose probabilities.

\noindent\textbf{Role in the proof.}
Use in this chapter. It is a reusable checked theorem cited by later proof scripts.

\noindent\textbf{Proof-script interpretation.}
Proof-script reading. The machine proof decomposes the theorem into rewrites, program-logic obligations, relational equivalences, and arithmetic side conditions.
\emph{Main tactics visible in the proof script:} \ec{rewrite}, \ec{smt}.
\emph{Named identifiers syntactically cited in the proof block:}
\begin{lstlisting}[language=EasyCrypt,basicstyle=\ttfamily\scriptsize]
bit_range
keccak_lane_bit
keccak_lane_of_int
nth_mkseq
\end{lstlisting}

\end{formaltheorem}

\begin{formaltheorem}[EasyCrypt line 374]
\noindent\textbf{EasyCrypt name.}
\begin{lstlisting}[language=EasyCrypt,basicstyle=\ttfamily\scriptsize]
keccak_state_lane_bitE
\end{lstlisting}
\noindent\textbf{Checked EasyCrypt statement.}
\begin{lstlisting}[language=EasyCrypt,basicstyle=\ttfamily\scriptsize]
lemma keccak_state_lane_bitE st x y bit :
  keccak_lane_bit (keccak_state_lane st x y) bit =
  keccak_lane_bit
    (nth keccak_lane_zero st (keccak_lane_index x y))
    bit.
\end{lstlisting}
\noindent\textbf{Formal mathematical reading.}
Mathematical theorem. This is an equality or definitional expansion used for rewriting and normalization.

\noindent\textbf{Role in the proof.}
Use in this chapter. It preserves an invariant across an oracle call, adversary call, or game procedure.

\noindent\textbf{Proof-script interpretation.}
No proof block was collected for this item.

\end{formaltheorem}

\begin{formaltheorem}[EasyCrypt line 402]
\noindent\textbf{EasyCrypt name.}
\begin{lstlisting}[language=EasyCrypt,basicstyle=\ttfamily\scriptsize]
keccak_theta_c_bitE
\end{lstlisting}
\noindent\textbf{Checked EasyCrypt statement.}
\begin{lstlisting}[language=EasyCrypt,basicstyle=\ttfamily\scriptsize]
lemma keccak_theta_c_bitE st x bit :
  0 <= bit < keccak_lane_bits =>
  keccak_lane_bit (keccak_theta_c st x) bit =
  ((((keccak_lane_bit (keccak_state_lane st x 0) bit <>
      keccak_lane_bit (keccak_state_lane st x 1) bit) <>
     keccak_lane_bit (keccak_state_lane st x 2) bit) <>
    keccak_lane_bit (keccak_state_lane st x 3) bit) <>
   keccak_lane_bit (keccak_state_lane st x 4) bit).
\end{lstlisting}
\noindent\textbf{Formal mathematical reading.}
Mathematical theorem. This is an inequality, usually an arithmetic loss bound, event inclusion bound, or monotonicity fact used to compose probabilities.

\noindent\textbf{Role in the proof.}
Use in this chapter. It is a reusable checked theorem cited by later proof scripts.

\noindent\textbf{Proof-script interpretation.}
Proof-script reading. The machine proof decomposes the theorem into rewrites, program-logic obligations, relational equivalences, and arithmetic side conditions.
\emph{Main tactics visible in the proof script:} \ec{rewrite}.
\emph{Named identifiers syntactically cited in the proof block:}
\begin{lstlisting}[language=EasyCrypt,basicstyle=\ttfamily\scriptsize]
bit_range
keccak_lane_xor_bitE
keccak_theta_c
\end{lstlisting}

\end{formaltheorem}

\begin{formaltheorem}[EasyCrypt line 415]
\noindent\textbf{EasyCrypt name.}
\begin{lstlisting}[language=EasyCrypt,basicstyle=\ttfamily\scriptsize]
keccak_theta_c_bit_indexE
\end{lstlisting}
\noindent\textbf{Checked EasyCrypt statement.}
\begin{lstlisting}[language=EasyCrypt,basicstyle=\ttfamily\scriptsize]
lemma keccak_theta_c_bit_indexE st x bit :
  0 <= bit < keccak_lane_bits =>
  keccak_lane_bit (keccak_theta_c st x) bit =
  keccak_theta_c_bit_index_value st x bit.
\end{lstlisting}
\noindent\textbf{Formal mathematical reading.}
Mathematical theorem. This is an inequality, usually an arithmetic loss bound, event inclusion bound, or monotonicity fact used to compose probabilities.

\noindent\textbf{Role in the proof.}
Use in this chapter. It is a reusable checked theorem cited by later proof scripts.

\noindent\textbf{Proof-script interpretation.}
Proof-script reading. The machine proof decomposes the theorem into rewrites, program-logic obligations, relational equivalences, and arithmetic side conditions.
\emph{Main tactics visible in the proof script:} \ec{rewrite}.
\emph{Named identifiers syntactically cited in the proof block:}
\begin{lstlisting}[language=EasyCrypt,basicstyle=\ttfamily\scriptsize]
bit_range
keccak_theta_c_bitE
keccak_theta_c_bit_index_value
\end{lstlisting}

\end{formaltheorem}

\begin{formaltheorem}[EasyCrypt line 425]
\noindent\textbf{EasyCrypt name.}
\begin{lstlisting}[language=EasyCrypt,basicstyle=\ttfamily\scriptsize]
keccak_theta_d_bitE
\end{lstlisting}
\noindent\textbf{Checked EasyCrypt statement.}
\begin{lstlisting}[language=EasyCrypt,basicstyle=\ttfamily\scriptsize]
lemma keccak_theta_d_bitE st x bit :
  0 <= bit < keccak_lane_bits =>
  keccak_lane_bit (keccak_theta_d st x) bit =
  (keccak_lane_bit (keccak_theta_c st (x + 4)) bit <>
   keccak_lane_bit
     (keccak_theta_c st (x + 1))
     ((bit - 1) %% keccak_lane_bits)).
\end{lstlisting}
\noindent\textbf{Formal mathematical reading.}
Mathematical theorem. This is an inequality, usually an arithmetic loss bound, event inclusion bound, or monotonicity fact used to compose probabilities.

\noindent\textbf{Role in the proof.}
Use in this chapter. It is a reusable checked theorem cited by later proof scripts.

\noindent\textbf{Proof-script interpretation.}
Proof-script reading. The machine proof decomposes the theorem into rewrites, program-logic obligations, relational equivalences, and arithmetic side conditions.
\emph{Main tactics visible in the proof script:} \ec{rewrite}, \ec{smt}.
\emph{Named identifiers syntactically cited in the proof block:}
\begin{lstlisting}[language=EasyCrypt,basicstyle=\ttfamily\scriptsize]
bit_range
keccak_lane_rotl_bitE
keccak_lane_xor_bitE
keccak_theta_d
\end{lstlisting}

\end{formaltheorem}

\begin{formaltheorem}[EasyCrypt line 439]
\noindent\textbf{EasyCrypt name.}
\begin{lstlisting}[language=EasyCrypt,basicstyle=\ttfamily\scriptsize]
keccak_theta_d_bit_indexE
\end{lstlisting}
\noindent\textbf{Checked EasyCrypt statement.}
\begin{lstlisting}[language=EasyCrypt,basicstyle=\ttfamily\scriptsize]
lemma keccak_theta_d_bit_indexE st x bit :
  0 <= bit < keccak_lane_bits =>
  keccak_lane_bit (keccak_theta_d st x) bit =
  keccak_theta_d_bit_index_value st x bit.
\end{lstlisting}
\noindent\textbf{Formal mathematical reading.}
Mathematical theorem. This is an inequality, usually an arithmetic loss bound, event inclusion bound, or monotonicity fact used to compose probabilities.

\noindent\textbf{Role in the proof.}
Use in this chapter. It is a reusable checked theorem cited by later proof scripts.

\noindent\textbf{Proof-script interpretation.}
Proof-script reading. The machine proof decomposes the theorem into rewrites, program-logic obligations, relational equivalences, and arithmetic side conditions.
\emph{Main tactics visible in the proof script:} \ec{rewrite}, \ec{smt}.
\emph{Named identifiers syntactically cited in the proof block:}
\begin{lstlisting}[language=EasyCrypt,basicstyle=\ttfamily\scriptsize]
bit_range
keccak_lane_bits
keccak_theta_c_bit_indexE
keccak_theta_d_bitE
keccak_theta_d_bit_index_value
\end{lstlisting}

\end{formaltheorem}

\begin{formaltheorem}[EasyCrypt line 452]
\noindent\textbf{EasyCrypt name.}
\begin{lstlisting}[language=EasyCrypt,basicstyle=\ttfamily\scriptsize]
keccak_theta_bitE
\end{lstlisting}
\noindent\textbf{Checked EasyCrypt statement.}
\begin{lstlisting}[language=EasyCrypt,basicstyle=\ttfamily\scriptsize]
lemma keccak_theta_bitE st lane bit :
  0 <= lane < keccak_state_lanes =>
  0 <= bit < keccak_lane_bits =>
  keccak_lane_bit
    (nth keccak_lane_zero (keccak_theta st) lane)
    bit =
  (let x = lane %% 5 in
   let y = lane %/ 5 in
   keccak_lane_bit (keccak_state_lane st x y) bit <>
   keccak_lane_bit (keccak_theta_d st x) bit).
\end{lstlisting}
\noindent\textbf{Formal mathematical reading.}
Mathematical theorem. This is an inequality, usually an arithmetic loss bound, event inclusion bound, or monotonicity fact used to compose probabilities.

\noindent\textbf{Role in the proof.}
Use in this chapter. It is a reusable checked theorem cited by later proof scripts.

\noindent\textbf{Proof-script interpretation.}
Proof-script reading. The machine proof decomposes the theorem into rewrites, program-logic obligations, relational equivalences, and arithmetic side conditions.
\emph{Main tactics visible in the proof script:} \ec{rewrite}, \ec{smt}.
\emph{Named identifiers syntactically cited in the proof block:}
\begin{lstlisting}[language=EasyCrypt,basicstyle=\ttfamily\scriptsize]
bit_range
keccak_lane_xor_bitE
keccak_theta
lane_range
nth_mkseq
\end{lstlisting}

\end{formaltheorem}

\begin{formaltheorem}[EasyCrypt line 469]
\noindent\textbf{EasyCrypt name.}
\begin{lstlisting}[language=EasyCrypt,basicstyle=\ttfamily\scriptsize]
keccak_theta_bit_indexE
\end{lstlisting}
\noindent\textbf{Checked EasyCrypt statement.}
\begin{lstlisting}[language=EasyCrypt,basicstyle=\ttfamily\scriptsize]
lemma keccak_theta_bit_indexE st lane bit :
  0 <= lane < keccak_state_lanes =>
  0 <= bit < keccak_lane_bits =>
  keccak_lane_bit
    (nth keccak_lane_zero (keccak_theta st) lane)
    bit =
  (let x = lane %% 5 in
   let y = lane %/ 5 in
   keccak_lane_bit (nth keccak_lane_zero st (keccak_lane_index x y)) bit <>
   keccak_lane_bit (keccak_theta_d st x) bit).
\end{lstlisting}
\noindent\textbf{Formal mathematical reading.}
Mathematical theorem. This is an inequality, usually an arithmetic loss bound, event inclusion bound, or monotonicity fact used to compose probabilities.

\noindent\textbf{Role in the proof.}
Use in this chapter. It is a reusable checked theorem cited by later proof scripts.

\noindent\textbf{Proof-script interpretation.}
Proof-script reading. The machine proof decomposes the theorem into rewrites, program-logic obligations, relational equivalences, and arithmetic side conditions.
\emph{Main tactics visible in the proof script:} \ec{rewrite}.
\emph{Named identifiers syntactically cited in the proof block:}
\begin{lstlisting}[language=EasyCrypt,basicstyle=\ttfamily\scriptsize]
bit_range
keccak_theta_bitE
lane_range
\end{lstlisting}

\end{formaltheorem}

\begin{formaltheorem}[EasyCrypt line 484]
\noindent\textbf{EasyCrypt name.}
\begin{lstlisting}[language=EasyCrypt,basicstyle=\ttfamily\scriptsize]
keccak_rho_pi_bitE
\end{lstlisting}
\noindent\textbf{Checked EasyCrypt statement.}
\begin{lstlisting}[language=EasyCrypt,basicstyle=\ttfamily\scriptsize]
lemma keccak_rho_pi_bitE st lane bit :
  0 <= lane < keccak_state_lanes =>
  0 <= bit < keccak_lane_bits =>
  keccak_lane_bit
    (nth keccak_lane_zero (keccak_rho_pi st) lane)
    bit =
  (let x = lane %% 5 in
   let y = lane %/ 5 in
   let sx = (x + 3 * y) %% 5 in
   let sy = x in
   keccak_lane_bit
     (keccak_state_lane st sx sy)
     ((bit - nth 0 keccak_rho_offsets (keccak_lane_index sx sy)) %%
      keccak_lane_bits)).
\end{lstlisting}
\noindent\textbf{Formal mathematical reading.}
Mathematical theorem. This is an inequality, usually an arithmetic loss bound, event inclusion bound, or monotonicity fact used to compose probabilities.

\noindent\textbf{Role in the proof.}
Use in this chapter. It is a reusable checked theorem cited by later proof scripts.

\noindent\textbf{Proof-script interpretation.}
Proof-script reading. The machine proof decomposes the theorem into rewrites, program-logic obligations, relational equivalences, and arithmetic side conditions.
\emph{Main tactics visible in the proof script:} \ec{rewrite}, \ec{smt}.
\emph{Named identifiers syntactically cited in the proof block:}
\begin{lstlisting}[language=EasyCrypt,basicstyle=\ttfamily\scriptsize]
bit_range
keccak_lane_rotl_bitE
keccak_rho_pi
lane_range
nth_mkseq
\end{lstlisting}

\end{formaltheorem}

\begin{formaltheorem}[EasyCrypt line 505]
\noindent\textbf{EasyCrypt name.}
\begin{lstlisting}[language=EasyCrypt,basicstyle=\ttfamily\scriptsize]
keccak_rho_pi_bit_indexE
\end{lstlisting}
\noindent\textbf{Checked EasyCrypt statement.}
\begin{lstlisting}[language=EasyCrypt,basicstyle=\ttfamily\scriptsize]
lemma keccak_rho_pi_bit_indexE st lane bit :
  0 <= lane < keccak_state_lanes =>
  0 <= bit < keccak_lane_bits =>
  keccak_lane_bit
    (nth keccak_lane_zero (keccak_rho_pi st) lane)
    bit =
  (let x = lane %% 5 in
   let y = lane %/ 5 in
   let sx = (x + 3 * y) %% 5 in
   let sy = x in
   keccak_lane_bit
     (nth keccak_lane_zero st (keccak_lane_index sx sy))
     ((bit - nth 0 keccak_rho_offsets (keccak_lane_index sx sy)) %%
      keccak_lane_bits)).
\end{lstlisting}
\noindent\textbf{Formal mathematical reading.}
Mathematical theorem. This is an inequality, usually an arithmetic loss bound, event inclusion bound, or monotonicity fact used to compose probabilities.

\noindent\textbf{Role in the proof.}
Use in this chapter. It is a reusable checked theorem cited by later proof scripts.

\noindent\textbf{Proof-script interpretation.}
Proof-script reading. The machine proof decomposes the theorem into rewrites, program-logic obligations, relational equivalences, and arithmetic side conditions.
\emph{Main tactics visible in the proof script:} \ec{rewrite}.
\emph{Named identifiers syntactically cited in the proof block:}
\begin{lstlisting}[language=EasyCrypt,basicstyle=\ttfamily\scriptsize]
bit_range
keccak_rho_pi_bitE
lane_range
\end{lstlisting}

\end{formaltheorem}

\begin{formaltheorem}[EasyCrypt line 524]
\noindent\textbf{EasyCrypt name.}
\begin{lstlisting}[language=EasyCrypt,basicstyle=\ttfamily\scriptsize]
keccak_chi_bitE
\end{lstlisting}
\noindent\textbf{Checked EasyCrypt statement.}
\begin{lstlisting}[language=EasyCrypt,basicstyle=\ttfamily\scriptsize]
lemma keccak_chi_bitE st lane bit :
  0 <= lane < keccak_state_lanes =>
  0 <= bit < keccak_lane_bits =>
  keccak_lane_bit
    (nth keccak_lane_zero (keccak_chi st) lane)
    bit =
  (let x = lane %% 5 in
   let y = lane %/ 5 in
   keccak_lane_bit (keccak_state_lane st x y) bit <>
   ((if keccak_lane_bit (keccak_state_lane st (x + 1) y) bit
     then false else true) /\
    keccak_lane_bit (keccak_state_lane st (x + 2) y) bit)).
\end{lstlisting}
\noindent\textbf{Formal mathematical reading.}
Mathematical theorem. This is an inequality, usually an arithmetic loss bound, event inclusion bound, or monotonicity fact used to compose probabilities.

\noindent\textbf{Role in the proof.}
Use in this chapter. It is a reusable checked theorem cited by later proof scripts.

\noindent\textbf{Proof-script interpretation.}
Proof-script reading. The machine proof decomposes the theorem into rewrites, program-logic obligations, relational equivalences, and arithmetic side conditions.
\emph{Main tactics visible in the proof script:} \ec{rewrite}, \ec{smt}.
\emph{Named identifiers syntactically cited in the proof block:}
\begin{lstlisting}[language=EasyCrypt,basicstyle=\ttfamily\scriptsize]
bit_range
keccak_chi
keccak_lane_and_bitE
keccak_lane_not_bitE
keccak_lane_xor_bitE
lane_range
nth_mkseq
\end{lstlisting}

\end{formaltheorem}

\begin{formaltheorem}[EasyCrypt line 545]
\noindent\textbf{EasyCrypt name.}
\begin{lstlisting}[language=EasyCrypt,basicstyle=\ttfamily\scriptsize]
keccak_chi_bit_indexE
\end{lstlisting}
\noindent\textbf{Checked EasyCrypt statement.}
\begin{lstlisting}[language=EasyCrypt,basicstyle=\ttfamily\scriptsize]
lemma keccak_chi_bit_indexE st lane bit :
  0 <= lane < keccak_state_lanes =>
  0 <= bit < keccak_lane_bits =>
  keccak_lane_bit
    (nth keccak_lane_zero (keccak_chi st) lane)
    bit =
  (let x = lane %% 5 in
   let y = lane %/ 5 in
   keccak_lane_bit (nth keccak_lane_zero st (keccak_lane_index x y)) bit <>
   ((if keccak_lane_bit
          (nth keccak_lane_zero st (keccak_lane_index (x + 1) y))
          bit
     then false else true) /\
    keccak_lane_bit
      (nth keccak_lane_zero st (keccak_lane_index (x + 2) y))
      bit)).
\end{lstlisting}
\noindent\textbf{Formal mathematical reading.}
Mathematical theorem. This is an inequality, usually an arithmetic loss bound, event inclusion bound, or monotonicity fact used to compose probabilities.

\noindent\textbf{Role in the proof.}
Use in this chapter. It is a reusable checked theorem cited by later proof scripts.

\noindent\textbf{Proof-script interpretation.}
Proof-script reading. The machine proof decomposes the theorem into rewrites, program-logic obligations, relational equivalences, and arithmetic side conditions.
\emph{Main tactics visible in the proof script:} \ec{rewrite}.
\emph{Named identifiers syntactically cited in the proof block:}
\begin{lstlisting}[language=EasyCrypt,basicstyle=\ttfamily\scriptsize]
bit_range
keccak_chi_bitE
lane_range
\end{lstlisting}

\end{formaltheorem}

\begin{formaltheorem}[EasyCrypt line 566]
\noindent\textbf{EasyCrypt name.}
\begin{lstlisting}[language=EasyCrypt,basicstyle=\ttfamily\scriptsize]
keccak_iota_bitE
\end{lstlisting}
\noindent\textbf{Checked EasyCrypt statement.}
\begin{lstlisting}[language=EasyCrypt,basicstyle=\ttfamily\scriptsize]
lemma keccak_iota_bitE st rc lane bit :
  0 <= lane < keccak_state_lanes =>
  0 <= bit < keccak_lane_bits =>
  keccak_lane_bit
    (nth keccak_lane_zero (keccak_iota st rc) lane)
    bit =
  if lane = 0
  then keccak_lane_bit (nth keccak_lane_zero st 0) bit <>
       keccak_lane_bit rc bit
  else keccak_lane_bit (nth keccak_lane_zero st lane) bit.
\end{lstlisting}
\noindent\textbf{Formal mathematical reading.}
Mathematical theorem. This is an inequality, usually an arithmetic loss bound, event inclusion bound, or monotonicity fact used to compose probabilities.

\noindent\textbf{Role in the proof.}
Use in this chapter. It is a reusable checked theorem cited by later proof scripts.

\noindent\textbf{Proof-script interpretation.}
Proof-script reading. The machine proof decomposes the theorem into rewrites, program-logic obligations, relational equivalences, and arithmetic side conditions.
\emph{Main tactics visible in the proof script:} \ec{rewrite}, \ec{smt}, \ec{case}.
\emph{Named identifiers syntactically cited in the proof block:}
\begin{lstlisting}[language=EasyCrypt,basicstyle=\ttfamily\scriptsize]
bit_range
keccak_iota
keccak_lane_xor_bitE
lane
lane0
lane_range
nth_mkseq
\end{lstlisting}

\end{formaltheorem}

\begin{formaltheorem}[EasyCrypt line 585]
\noindent\textbf{EasyCrypt name.}
\begin{lstlisting}[language=EasyCrypt,basicstyle=\ttfamily\scriptsize]
keccak_round_bitE
\end{lstlisting}
\noindent\textbf{Checked EasyCrypt statement.}
\begin{lstlisting}[language=EasyCrypt,basicstyle=\ttfamily\scriptsize]
lemma keccak_round_bitE st rc lane bit :
  0 <= lane < keccak_state_lanes =>
  0 <= bit < keccak_lane_bits =>
  keccak_lane_bit
    (nth keccak_lane_zero (keccak_round st rc) lane)
    bit =
  if lane = 0
  then
    keccak_lane_bit
      (nth keccak_lane_zero
        (keccak_chi (keccak_rho_pi (keccak_theta st)))
        0)
      bit <>
    keccak_lane_bit rc bit
  else
    keccak_lane_bit
      (nth keccak_lane_zero
        (keccak_chi (keccak_rho_pi (keccak_theta st)))
        lane)
      bit.
\end{lstlisting}
\noindent\textbf{Formal mathematical reading.}
Mathematical theorem. This is an inequality, usually an arithmetic loss bound, event inclusion bound, or monotonicity fact used to compose probabilities.

\noindent\textbf{Role in the proof.}
Use in this chapter. It is a reusable checked theorem cited by later proof scripts.

\noindent\textbf{Proof-script interpretation.}
Proof-script reading. The machine proof decomposes the theorem into rewrites, program-logic obligations, relational equivalences, and arithmetic side conditions.
\emph{Main tactics visible in the proof script:} \ec{rewrite}.
\emph{Named identifiers syntactically cited in the proof block:}
\begin{lstlisting}[language=EasyCrypt,basicstyle=\ttfamily\scriptsize]
bit_range
keccak_iota_bitE
keccak_round
lane_range
\end{lstlisting}

\end{formaltheorem}

\begin{formaltheorem}[EasyCrypt line 611]
\noindent\textbf{EasyCrypt name.}
\begin{lstlisting}[language=EasyCrypt,basicstyle=\ttfamily\scriptsize]
keccak_bytes_of_lanes_size
\end{lstlisting}
\noindent\textbf{Checked EasyCrypt statement.}
\begin{lstlisting}[language=EasyCrypt,basicstyle=\ttfamily\scriptsize]
lemma keccak_bytes_of_lanes_size st :
  size (keccak_bytes_of_lanes st) = keccak_state_bytes.
\end{lstlisting}
\noindent\textbf{Formal mathematical reading.}
Mathematical theorem. This is an equality or definitional expansion used for rewriting and normalization.

\noindent\textbf{Role in the proof.}
Use in this chapter. It is a reusable checked theorem cited by later proof scripts.

\noindent\textbf{Proof-script interpretation.}
No proof block was collected for this item.

\end{formaltheorem}

\begin{formaltheorem}[EasyCrypt line 615]
\noindent\textbf{EasyCrypt name.}
\begin{lstlisting}[language=EasyCrypt,basicstyle=\ttfamily\scriptsize]
keccak_bytes_of_lanes_byteE
\end{lstlisting}
\noindent\textbf{Checked EasyCrypt statement.}
\begin{lstlisting}[language=EasyCrypt,basicstyle=\ttfamily\scriptsize]
lemma keccak_bytes_of_lanes_byteE st i :
  0 <= i < keccak_state_bytes =>
  nth 0 (keccak_bytes_of_lanes st) i =
  keccak_lane_to_byte
    (nth keccak_lane_zero st (i %/ 8))
    (i %% 8).
\end{lstlisting}
\noindent\textbf{Formal mathematical reading.}
Mathematical theorem. This is an inequality, usually an arithmetic loss bound, event inclusion bound, or monotonicity fact used to compose probabilities.

\noindent\textbf{Role in the proof.}
Use in this chapter. It is a reusable checked theorem cited by later proof scripts.

\noindent\textbf{Proof-script interpretation.}
Proof-script reading. The machine proof decomposes the theorem into rewrites, program-logic obligations, relational equivalences, and arithmetic side conditions.
\emph{Main tactics visible in the proof script:} \ec{rewrite}.
\emph{Named identifiers syntactically cited in the proof block:}
\begin{lstlisting}[language=EasyCrypt,basicstyle=\ttfamily\scriptsize]
i_range
keccak_bytes_of_lanes
nth_mkseq
\end{lstlisting}

\end{formaltheorem}

\begin{formaltheorem}[EasyCrypt line 626]
\noindent\textbf{EasyCrypt name.}
\begin{lstlisting}[language=EasyCrypt,basicstyle=\ttfamily\scriptsize]
keccak_lane_to_byte_bitsE
\end{lstlisting}
\noindent\textbf{Checked EasyCrypt statement.}
\begin{lstlisting}[language=EasyCrypt,basicstyle=\ttfamily\scriptsize]
lemma keccak_lane_to_byte_bitsE lane byte_index :
  keccak_lane_to_byte lane byte_index =
  (if keccak_lane_bit lane (8 * byte_index) then 2 ^ 0 else 0) +
  ((if keccak_lane_bit lane (8 * byte_index + 1) then 2 ^ 1 else 0) +
  ((if keccak_lane_bit lane (8 * byte_index + 2) then 2 ^ 2 else 0) +
  ((if keccak_lane_bit lane (8 * byte_index + 3) then 2 ^ 3 else 0) +
  ((if keccak_lane_bit lane (8 * byte_index + 4) then 2 ^ 4 else 0) +
  ((if keccak_lane_bit lane (8 * byte_index + 5) then 2 ^ 5 else 0) +
  ((if keccak_lane_bit lane (8 * byte_index + 6) then 2 ^ 6 else 0) +
   (if keccak_lane_bit lane (8 * byte_index + 7) then 2 ^ 7 else 0))))))).
\end{lstlisting}
\noindent\textbf{Formal mathematical reading.}
Mathematical theorem. This is an equality or definitional expansion used for rewriting and normalization.

\noindent\textbf{Role in the proof.}
Use in this chapter. It is a reusable checked theorem cited by later proof scripts.

\noindent\textbf{Proof-script interpretation.}
Proof-script reading. The machine proof decomposes the theorem into rewrites, program-logic obligations, relational equivalences, and arithmetic side conditions.
\emph{Main tactics visible in the proof script:} \ec{rewrite}.
\emph{Named identifiers syntactically cited in the proof block:}
\begin{lstlisting}[language=EasyCrypt,basicstyle=\ttfamily\scriptsize]
iota0
iotaS
keccak_lane_to_byte
range
sumz
\end{lstlisting}

\end{formaltheorem}

\begin{formaltheorem}[EasyCrypt line 650]
\noindent\textbf{EasyCrypt name.}
\begin{lstlisting}[language=EasyCrypt,basicstyle=\ttfamily\scriptsize]
keccak_lane_to_byte_from_bitsE
\end{lstlisting}
\noindent\textbf{Checked EasyCrypt statement.}
\begin{lstlisting}[language=EasyCrypt,basicstyle=\ttfamily\scriptsize]
lemma keccak_lane_to_byte_from_bitsE
  lane byte_index b0 b1 b2 b3 b4 b5 b6 b7 :
  keccak_lane_bit lane (8 * byte_index) = b0 =>
  keccak_lane_bit lane (8 * byte_index + 1) = b1 =>
  keccak_lane_bit lane (8 * byte_index + 2) = b2 =>
  keccak_lane_bit lane (8 * byte_index + 3) = b3 =>
  keccak_lane_bit lane (8 * byte_index + 4) = b4 =>
  keccak_lane_bit lane (8 * byte_index + 5) = b5 =>
  keccak_lane_bit lane (8 * byte_index + 6) = b6 =>
  keccak_lane_bit lane (8 * byte_index + 7) = b7 =>
  keccak_lane_to_byte lane byte_index =
  (if b0 then 1 else 0) +
  ((if b1 then 2 else 0) +
  ((if b2 then 4 else 0) +
  ((if b3 then 8 else 0) +
  ((if b4 then 16 else 0) +
  ((if b5 then 32 else 0) +
  ((if b6 then 64 else 0) +
   (if b7 then 128 else 0))))))).
\end{lstlisting}
\noindent\textbf{Formal mathematical reading.}
Mathematical theorem. This is an implication, normally turning one invariant, validity predicate, or proof obligation into another.

\noindent\textbf{Role in the proof.}
Use in this chapter. It is a reusable checked theorem cited by later proof scripts.

\noindent\textbf{Proof-script interpretation.}
Proof-script reading. The machine proof decomposes the theorem into rewrites, program-logic obligations, relational equivalences, and arithmetic side conditions.
\emph{Main tactics visible in the proof script:} \ec{rewrite}.
\emph{Named identifiers syntactically cited in the proof block:}
\begin{lstlisting}[language=EasyCrypt,basicstyle=\ttfamily\scriptsize]
h0
h1
h2
h3
h4
h5
h6
h7
keccak_lane_to_byte_bitsE
keccak_pow2_0E
keccak_pow2_1E
keccak_pow2_2E
keccak_pow2_3E
keccak_pow2_4E
keccak_pow2_5E
keccak_pow2_6E
keccak_pow2_7E
\end{lstlisting}

\end{formaltheorem}

\begin{formaltheorem}[EasyCrypt line 676]
\noindent\textbf{EasyCrypt name.}
\begin{lstlisting}[language=EasyCrypt,basicstyle=\ttfamily\scriptsize]
keccak_lane_to_byte_70E
\end{lstlisting}
\noindent\textbf{Checked EasyCrypt statement.}
\begin{lstlisting}[language=EasyCrypt,basicstyle=\ttfamily\scriptsize]
lemma keccak_lane_to_byte_70E lane byte_index :
  keccak_lane_bit lane (8 * byte_index) = false =>
  keccak_lane_bit lane (8 * byte_index + 1) = true =>
  keccak_lane_bit lane (8 * byte_index + 2) = true =>
  keccak_lane_bit lane (8 * byte_index + 3) = false =>
  keccak_lane_bit lane (8 * byte_index + 4) = false =>
  keccak_lane_bit lane (8 * byte_index + 5) = false =>
  keccak_lane_bit lane (8 * byte_index + 6) = true =>
  keccak_lane_bit lane (8 * byte_index + 7) = false =>
  keccak_lane_to_byte lane byte_index = 70.
\end{lstlisting}
\noindent\textbf{Formal mathematical reading.}
Mathematical theorem. This is an implication, normally turning one invariant, validity predicate, or proof obligation into another.

\noindent\textbf{Role in the proof.}
Use in this chapter. It is a reusable checked theorem cited by later proof scripts.

\noindent\textbf{Proof-script interpretation.}
Proof-script reading. The machine proof decomposes the theorem into rewrites, program-logic obligations, relational equivalences, and arithmetic side conditions.
\emph{Main tactics visible in the proof script:} \ec{rewrite}, \ec{smt}.
\emph{Named identifiers syntactically cited in the proof block:}
\begin{lstlisting}[language=EasyCrypt,basicstyle=\ttfamily\scriptsize]
b0
b1
b2
b3
b4
b5
b6
b7
keccak_lane_to_byte_bitsE
keccak_pow2_1E
keccak_pow2_2E
keccak_pow2_6E
\end{lstlisting}

\end{formaltheorem}

\begin{formaltheorem}[EasyCrypt line 693]
\noindent\textbf{EasyCrypt name.}
\begin{lstlisting}[language=EasyCrypt,basicstyle=\ttfamily\scriptsize]
keccak_f1600_bytes_size
\end{lstlisting}
\noindent\textbf{Checked EasyCrypt statement.}
\begin{lstlisting}[language=EasyCrypt,basicstyle=\ttfamily\scriptsize]
lemma keccak_f1600_bytes_size st :
  size (keccak_f1600_bytes st) = keccak_state_bytes.
\end{lstlisting}
\noindent\textbf{Formal mathematical reading.}
Mathematical theorem. This is an equality or definitional expansion used for rewriting and normalization.

\noindent\textbf{Role in the proof.}
Use in this chapter. It is a reusable checked theorem cited by later proof scripts.

\noindent\textbf{Proof-script interpretation.}
No proof block was collected for this item.

\end{formaltheorem}

\medskip
\noindent\emph{End of generated formal inventory for \file{HAETAE_Keccak1600.ec}.}
